% ============================================================
% G1 Tank Robot -- Complete Laboratory Report
% Yahboom G1 Tank | Arduino UNO Version
% ============================================================
\documentclass[12pt,a4paper]{article}

% --- Encoding & fonts ---
\usepackage[T1]{fontenc}
\usepackage[utf8]{inputenc}
\usepackage{mathpazo}
\usepackage{helvet}

% --- Page layout ---
\usepackage[
  a4paper,
  top=25mm, bottom=25mm,
  left=25mm, right=25mm,
  headheight=15pt
]{geometry}

% --- Math ---
\usepackage{amsmath}
\usepackage{amssymb}
\usepackage{mathtools}
\usepackage{bm}

% --- Graphics & color ---
\usepackage{xcolor}
\usepackage{tikz}
\usetikzlibrary{shapes.geometric, arrows.meta, positioning, fit, backgrounds, calc}
\usepackage{graphicx}

% --- Tables ---
\usepackage{booktabs}
\usepackage{array}
\usepackage{multirow}
\usepackage{longtable}
\usepackage{tabularx}
\usepackage{colortbl}

% --- Code listings ---
\usepackage{listings}

% --- Boxes ---
\usepackage{tcolorbox}
\tcbuselibrary{breakable, skins, theorems}

% --- Headers/footers ---
\usepackage{fancyhdr}

% --- Section formatting ---
\usepackage{titlesec}

% --- Lists ---
\usepackage{enumitem}

% --- Hyperlinks ---
\usepackage{hyperref}

% --- Misc ---
\usepackage{mdframed}
\usepackage{microtype}
\usepackage{parskip}

% ============================================================
%  COLOR PALETTE
% ============================================================
\definecolor{PrimaryBlue}{RGB}{31,78,121}
\definecolor{MidBlue}{RGB}{46,117,182}
\definecolor{LightBlue}{RGB}{189,215,238}
\definecolor{VeryLightBlue}{RGB}{235,243,251}
\definecolor{DarkGray}{RGB}{64,64,64}
\definecolor{MedGray}{RGB}{120,120,120}
\definecolor{LightGray}{RGB}{242,242,242}
\definecolor{SuccessGreen}{RGB}{55,86,35}
\definecolor{SuccessBG}{RGB}{226,239,218}
\definecolor{WarnOrange}{RGB}{127,79,0}
\definecolor{WarnBG}{RGB}{255,242,204}
\definecolor{DangerRed}{RGB}{123,0,0}
\definecolor{DangerBG}{RGB}{252,228,214}
\definecolor{InfoBG}{RGB}{222,235,247}
\definecolor{CodeBG}{RGB}{30,30,30}
\definecolor{CodeText}{RGB}{212,212,212}
\definecolor{CodeKeyword}{RGB}{86,156,214}
\definecolor{CodeComment}{RGB}{106,153,85}
\definecolor{CodeString}{RGB}{206,145,120}
\definecolor{CodeNumber}{RGB}{181,206,168}
\definecolor{RuleLine}{RGB}{46,117,182}

% ============================================================
%  LISTINGS STYLE  (dark theme -- Arduino/C/Python)
% ============================================================
\lstdefinestyle{arduino}{
  language=C++,
  backgroundcolor=\color{CodeBG},
  basicstyle=\ttfamily\footnotesize\color{CodeText},
  keywordstyle=\color{CodeKeyword}\bfseries,
  commentstyle=\color{CodeComment}\itshape,
  stringstyle=\color{CodeString},
  numberstyle=\tiny\color{MedGray},
  numbers=left,
  numbersep=10pt,
  stepnumber=1,
  frame=single,
  framerule=0pt,
  rulecolor=\color{CodeBG},
  tabsize=2,
  breaklines=true,
  breakatwhitespace=false,
  showspaces=false,
  showstringspaces=false,
  captionpos=b,
  xleftmargin=16pt,
  xrightmargin=4pt,
  literate=
    {°}{{\textdegree}}1
    {Σ}{{$\Sigma$}}1
    {→}{{$\rightarrow$}}1,
}

\lstdefinestyle{python}{
  language=Python,
  backgroundcolor=\color{CodeBG},
  basicstyle=\ttfamily\footnotesize\color{CodeText},
  keywordstyle=\color{CodeKeyword}\bfseries,
  commentstyle=\color{CodeComment}\itshape,
  stringstyle=\color{CodeString},
  numberstyle=\tiny\color{MedGray},
  numbers=left,
  numbersep=10pt,
  stepnumber=1,
  frame=single,
  framerule=0pt,
  rulecolor=\color{CodeBG},
  tabsize=4,
  breaklines=true,
  showspaces=false,
  showstringspaces=false,
  captionpos=b,
  xleftmargin=16pt,
  xrightmargin=4pt,
}

% ============================================================
%  TCOLORBOX ENVIRONMENTS
% ============================================================
% Concept box (blue)
\newtcolorbox{conceptbox}[1]{
  colback=VeryLightBlue,
  colframe=MidBlue,
  coltitle=white,
  fonttitle=\bfseries\small,
  title={#1},
  breakable,
  boxrule=0.8pt,
  arc=4pt,
  left=8pt, right=8pt, top=6pt, bottom=6pt,
}

% Success / done box (green)
\newtcolorbox{donebox}[1]{
  colback=SuccessBG,
  colframe=SuccessGreen,
  coltitle=white,
  fonttitle=\bfseries\small,
  title={#1},
  breakable,
  boxrule=0.8pt,
  arc=4pt,
  left=8pt, right=8pt, top=6pt, bottom=6pt,
}

% Warning box (orange)
\newtcolorbox{warnbox}[1]{
  colback=WarnBG,
  colframe=WarnOrange,
  coltitle=white,
  fonttitle=\bfseries\small,
  title={#1},
  breakable,
  boxrule=0.8pt,
  arc=4pt,
  left=8pt, right=8pt, top=6pt, bottom=6pt,
}

% Danger / blocked box (red)
\newtcolorbox{dangerbox}[1]{
  colback=DangerBG,
  colframe=DangerRed,
  coltitle=white,
  fonttitle=\bfseries\small,
  title={#1},
  breakable,
  boxrule=0.8pt,
  arc=4pt,
  left=8pt, right=8pt, top=6pt, bottom=6pt,
}

% Math derivation box
\newtcolorbox{mathbox}[1]{
  colback=VeryLightBlue,
  colframe=PrimaryBlue,
  coltitle=white,
  fonttitle=\bfseries\small,
  title={#1},
  breakable,
  boxrule=1pt,
  arc=4pt,
  left=10pt, right=10pt, top=8pt, bottom=8pt,
}

% Key result box
\newtcolorbox{resultbox}{
  colback=LightBlue!40,
  colframe=PrimaryBlue,
  boxrule=1.5pt,
  arc=4pt,
  left=10pt, right=10pt, top=8pt, bottom=8pt,
}

% ============================================================
%  SECTION FORMATTING
% ============================================================
\titleformat{\section}
  {\color{PrimaryBlue}\Large\bfseries}
  {\thesection}{1em}{}
  [\color{MidBlue}\titlerule]

\titleformat{\subsection}
  {\color{MidBlue}\large\bfseries}
  {\thesubsection}{1em}{}

\titleformat{\subsubsection}
  {\color{DarkGray}\normalsize\bfseries}
  {\thesubsubsection}{1em}{}

\titlespacing*{\section}{0pt}{18pt}{8pt}
\titlespacing*{\subsection}{0pt}{14pt}{6pt}
\titlespacing*{\subsubsection}{0pt}{10pt}{4pt}

% ============================================================
%  HEADER / FOOTER
% ============================================================
\pagestyle{fancy}
\fancyhf{}
\fancyhead[L]{\small\color{MidBlue}\textbf{Yahboom G1 Tank Robot} --- Laboratory Report}
\fancyhead[R]{\small\color{MedGray}Arduino UNO Version}
\fancyfoot[L]{\small\color{MedGray}Shenzhen Yahboom Technology Co., Ltd.}
\fancyfoot[R]{\small\color{MedGray}Page \thepage}
\renewcommand{\headrulewidth}{0.5pt}
\renewcommand{\headrule}{\color{MidBlue}\hrule width\headwidth height\headrulewidth}
\renewcommand{\footrulewidth}{0.3pt}

% ============================================================
%  HYPERREF SETUP
% ============================================================
\hypersetup{
  colorlinks=true,
  linkcolor=MidBlue,
  urlcolor=MidBlue,
  citecolor=MidBlue,
  pdftitle={Yahboom G1 Tank Robot -- Laboratory Report},
  pdfauthor={Lab Student},
}

% ============================================================
%  CUSTOM COMMANDS
% ============================================================
\newcommand{\status}[2]{%
  \begin{tcolorbox}[
    colback=#1!15,colframe=#1,boxrule=0.6pt,arc=3pt,
    left=6pt,right=6pt,top=4pt,bottom=4pt
  ]
  \small\textbf{#2}
  \end{tcolorbox}%
}

% Shorthand for skew-symmetric matrix notation
\renewcommand{\skew}[1]{\left[#1\right]}
\newcommand{\R}{\mathbb{R}}
\newcommand{\SO}{\mathrm{SO}(3)}
\newcommand{\SE}{\mathrm{SE}(3)}
\newcommand{\vect}[1]{\bm{#1}}
\newcommand{\mat}[1]{\mathbf{#1}}
\newcommand{\norm}[1]{\left\|#1\right\|}

% ============================================================
%  TABLE COLUMN TYPES
% ============================================================
\newcolumntype{L}[1]{>{\raggedright\arraybackslash}p{#1}}
\newcolumntype{C}[1]{>{\centering\arraybackslash}p{#1}}
\newcolumntype{R}[1]{>{\raggedleft\arraybackslash}p{#1}}

% ============================================================
\begin{document}
% ============================================================

% ---- TITLE PAGE ----
\begin{titlepage}
  \centering
  \vspace*{2cm}

  {\color{PrimaryBlue}\rule{\linewidth}{2pt}}
  \vspace{0.5cm}

  {\color{PrimaryBlue}\Huge\bfseries Yahboom G1 Tank Robot\\[0.4em]
   \color{MidBlue}\Large Laboratory Tasks --- Complete Technical Report}

  \vspace{0.5cm}
  {\color{PrimaryBlue}\rule{\linewidth}{2pt}}

  \vspace{1.5cm}

  \begin{tcolorbox}[
    colback=VeryLightBlue,colframe=MidBlue,
    boxrule=1pt,arc=6pt,
    left=16pt,right=16pt,top=12pt,bottom=12pt,
    width=0.85\linewidth
  ]
  \centering
  \begin{tabular}{ll}
    \textbf{Platform:}    & Yahboom G1 Tank Robot \\[4pt]
    \textbf{Controller:}  & Arduino UNO (Bluetooth) \\[4pt]
    \textbf{Constraint:}  & \color{DangerRed}\bfseries Raspberry Pi image failed --- UNO only \\[4pt]
    \textbf{Tasks:}       & DOF, End-effector angles, POE, Per-actuator verification \\
  \end{tabular}
  \end{tcolorbox}

  \vspace{1.5cm}

  {\large\bfseries\color{PrimaryBlue} Tasks Covered}\\[0.6em]
  \begin{tabular}{cl}
    \rowcolor{LightBlue!50}
    \textbf{Task 1} & Degrees of Freedom (DOF) --- Grubler's Formula \\
    \textbf{Task 2} & Programming End Effectors at Specific Angles \\
    \rowcolor{LightBlue!50}
    \textbf{Task 3} & Verifying Positions using Product of Exponentials (POE) \\
    \textbf{Task 4} & Per-Actuator Position Programming \& Angle Verification \\
  \end{tabular}

  \vfill

  \begin{tcolorbox}[
    colback=DangerBG,colframe=DangerRed,
    boxrule=0.8pt,arc=4pt,
    left=10pt,right=10pt,top=6pt,bottom=6pt,
    width=0.85\linewidth
  ]
  \centering\small
  \textbf{Hardware Constraint Note:} The Yahboom Raspberry Pi OS image failed to boot on the
  Pi board during the lab session. The Arduino UNO version was used instead.
  Camera pan and tilt servos (DOF 4 \& 5) are wired to Pi GPIO only and could
  not be physically tested. All analysis for these is marked \textbf{THEORETICAL}.
  \end{tcolorbox}

  \vspace{0.5cm}
  {\small\color{MedGray} Shenzhen Yahboom Technology Co., Ltd. \textbar\
   \texttt{www.yahboom.com}}
\end{titlepage}

% ---- TABLE OF CONTENTS ----
\tableofcontents
\newpage

% ============================================================
\section{Hardware Constraint \& Setup}
% ============================================================

\begin{dangerbox}{CONSTRAINT --- Raspberry Pi Image Failure}
During the laboratory session, the Yahboom OS image provided for the Raspberry Pi
board failed to boot after multiple attempts. The decision was made to switch to
the Arduino UNO configuration. The Yahboom Bluetooth app then successfully
connected, enabling movement and ultrasonic servo control. However, the camera
pan and tilt servos are wired exclusively to the Raspberry Pi GPIO header and
are inaccessible from the UNO.
\end{dangerbox}

\vspace{6pt}
\begin{table}[h!]
\centering
\caption{Component operational status}
\renewcommand{\arraystretch}{1.3}
\begin{tabularx}{\linewidth}{L{3.5cm} L{2.8cm} C{2.5cm} X}
\toprule
\rowcolor{MidBlue}
\color{white}\textbf{Component} &
\color{white}\textbf{Controller} &
\color{white}\textbf{Status} &
\color{white}\textbf{Impact} \\
\midrule
Left \& Right Track Motors & Arduino UNO & \cellcolor{SuccessBG}\color{SuccessGreen}OPERATIONAL & Full physical testing done \\
Ultrasonic Servo (J3) & UNO Pin 3 & \cellcolor{SuccessBG}\color{SuccessGreen}OPERATIONAL & Full physical testing done \\
HC-SR04 Sensor & UNO Pins 12/13 & \cellcolor{SuccessBG}\color{SuccessGreen}OPERATIONAL & Distance readings verified \\
Camera Pan Servo (J4) & Pi GPIO J2 & \cellcolor{DangerBG}\color{DangerRed}INACCESSIBLE & Theoretical analysis only \\
Camera Tilt Servo (J5) & Pi GPIO J2 & \cellcolor{DangerBG}\color{DangerRed}INACCESSIBLE & Theoretical analysis only \\
Pi Camera (CSI) & Raspberry Pi & \cellcolor{DangerBG}\color{DangerRed}INACCESSIBLE & No visual confirmation \\
RGB LED & UNO Pins 9/10/11 & \cellcolor{SuccessBG}\color{SuccessGreen}OPERATIONAL & Used for visual feedback \\
\bottomrule
\end{tabularx}
\end{table}

% ============================================================
\section{Task 1 --- Degrees of Freedom (DOF) Analysis}
% ============================================================

\begin{donebox}{STATUS: FULLY COMPLETE --- No hardware required}
DOF analysis is purely analytical. All results physically verified for the 3
accessible actuators; DOF 4 \& 5 verified mathematically.
\end{donebox}

% ------ 1.1 ------
\subsection{Fundamental Concepts}

\begin{conceptbox}{What is DOF?}
\textbf{Degree of Freedom (DOF)} is the number of independent parameters needed
to completely specify the configuration of a mechanical system. It answers the
question: \emph{``How many things can move independently?''}

A system with $n$ DOF requires exactly $n$ independent coordinates to describe
its configuration at any instant.
\end{conceptbox}

\vspace{6pt}

\begin{table}[h!]
\centering
\caption{Key concepts in mechanism analysis}
\renewcommand{\arraystretch}{1.3}
\begin{tabularx}{\linewidth}{L{3cm} X}
\toprule
\rowcolor{LightBlue}
\textbf{Concept} & \textbf{Definition} \\
\midrule
\textbf{Link} & Any rigid body that does not deform during motion. Bodies bolted
  together with no relative motion count as \emph{one} link. \\
\midrule
\textbf{Joint} & A connection between two links that constrains their relative
  motion. Each joint allows a specific number of relative DOF. \\
\midrule
\textbf{Ground Link} & The fixed reference body. It is assigned \emph{no joint
  relative to itself}. Everything else moves relative to it. This is the most
  commonly violated rule in DOF analysis. \\
\midrule
\textbf{Revolute Joint} & Allows exactly one rotation about a fixed axis.
  $f_i = 1$. All G1 Tank joints are revolute. \\
\bottomrule
\end{tabularx}
\end{table}


\textbf{Choosing $m = 3$ (planar):} The G1 Tank operates on a flat 2D ground
surface. The chassis does not fly, pitch, or roll. All mechanism motion is
analysed relative to the chassis (ground link), and the robot's operational
space is fundamentally planar. Therefore $m = 3$.

% ------ 1.3 ------
\subsection{Link Identification}

The chassis frame and the top plate (which holds the microcontroller) are
rigidly bolted together with no relative motion between them. By the
\emph{rigid body merging rule}, they are treated as one single link --- the
\textbf{ground link} $L_1$.

\begin{table}[h!]
\centering
\caption{All six links of the G1 Tank}
\renewcommand{\arraystretch}{1.3}
\begin{tabularx}{\linewidth}{C{1cm} L{3.5cm} X L{2.5cm}}
\toprule
\rowcolor{MidBlue}
\color{white}\textbf{Link} &
\color{white}\textbf{Name} &
\color{white}\textbf{Description} &
\color{white}\textbf{Type} \\
\midrule
$L_1$ & Chassis + Body & Bottom frame + top plate merged --- fixed reference & \cellcolor{WarnBG}\textbf{Ground link} \\
$L_2$ & Left Track & Left rubber track + sprocket, driven by left 370 DC motor & Moving \\
$L_3$ & Right Track & Right rubber track + sprocket, driven by right 370 DC motor & Moving \\
$L_4$ & Ultrasonic Sensor & HC-SR04 on servo bracket; pans left/right & Moving \\
$L_5$ & Camera Pan Bracket & Acrylic bracket on pan servo horn; rotates horizontally & Moving \\
$L_6$ & Camera + Tilt & Pi camera on tilt servo horn; rotates vertically on $L_5$ & Moving \\
\bottomrule
\end{tabularx}
\end{table}

\[
  \boxed{n = 6}
\]

% ------ 1.4 ------
\subsection{Joint Identification}

All joints are \textbf{revolute} ($f_i = 1$ each).

\begin{table}[h!]
\centering
\caption{All five joints of the G1 Tank}
\renewcommand{\arraystretch}{1.3}
\begin{tabularx}{\linewidth}{C{1cm} C{2cm} C{1.5cm} C{0.8cm} L{3cm} L{4.5cm}}
\toprule
\rowcolor{MidBlue}
\color{white}\textbf{Joint} &
\color{white}\textbf{Between} &
\color{white}\textbf{Type} &
\color{white}$f_i$ &
\color{white}\textbf{Actuator} &
\color{white}\textbf{Physical Description} \\
\midrule
$J_1$ & $L_1 \leftrightarrow L_2$ & Revolute & 1 & Left DC Motor & Left track rotates about horizontal axle \\
$J_2$ & $L_1 \leftrightarrow L_3$ & Revolute & 1 & Right DC Motor & Right track rotates about horizontal axle \\
$J_3$ & $L_1 \leftrightarrow L_4$ & Revolute & 1 & Ultrasonic Servo & Sensor pans left/right about vertical axis \\
$J_4$ & $L_1 \leftrightarrow L_5$ & Revolute & 1 & Camera Pan Servo & Pan bracket rotates horizontally \\
$J_5$ & $L_5 \leftrightarrow L_6$ & Revolute & 1 & Camera Tilt Servo & Camera tilts up/down about horizontal axis \\
\bottomrule
\end{tabularx}
\end{table}

\[
  \boxed{j = 5}, \qquad \boxed{\sum f_i = 5 \times 1 = 5}
\]

% ------ 1.5 ------
\subsection{Grubler Calculation}

\begin{mathbox}{G1 Tank --- DOF Calculation}
\begin{align*}
  \mathrm{DOF} &= m(n - 1 - j) + \sum f_i \\[6pt]
               &= 3(6 - 1 - 5) + 5 \\[6pt]
               &= 3(0) + 5 \\[6pt]
               &= \boxed{5}
\end{align*}
\end{mathbox}

% ============================================================
\newpage
\section{Task 2 --- Programming End Effectors at Specific Angles}
% ============================================================

\begin{warnbox}{STATUS: PARTIAL --- Ultrasonic + tracks physical; camera servos theoretical}
The ultrasonic servo (DOF 3) and both DC track motors (DOF 1 \& 2) are fully
accessible from the Arduino UNO and were physically programmed and verified.
Camera pan and tilt servos require the Raspberry Pi and could not be physically
commanded.
\end{warnbox}

% ------ 2.1 Background ------
\subsection{Background Knowledge --- How Servo Motors Work}

\begin{conceptbox}{What You Need to Know Before Programming a Servo}
A servo motor is a motor that can be commanded to move to a \textbf{specific
angular position} and hold it. It does this by using internal feedback (a
potentiometer or encoder). The controller speaks to the servo using
\textbf{PWM --- Pulse Width Modulation}:

\begin{itemize}[leftmargin=1.5em]
  \item A square wave signal is sent at \textbf{50\,Hz} (one pulse every 20\,ms).
  \item The \textbf{width of the HIGH pulse} tells the servo what angle to go to.
  \item Standard servos: 500\,$\mu$s pulse $= 0^\circ$, 2500\,$\mu$s pulse $= 180^\circ$.
  \item The servo \emph{ignores} the LOW part of the wave --- only the HIGH duration matters.
\end{itemize}

On the Arduino UNO, because the available hardware PWM pins are used for the
motors, the servo signal is generated in \textbf{software} (bit-banging) by the
existing code. This is slightly less precise than hardware PWM but sufficient for
a standard SG90-class servo.
\end{conceptbox}

\subsubsection{Pulse Width Formula (from existing code)}

The code uses:
\[
  \text{PulseWidth} [\mu\text{s}] = (\theta \times 11) + 500
\]

where $\theta$ is the angle in degrees. This maps:

\begin{center}
\renewcommand{\arraystretch}{1.2}
\begin{tabular}{ccc}
\toprule
\textbf{Angle} $\theta$ & \textbf{Pulse Width} & \textbf{Position} \\
\midrule
$0^\circ$   & $500\,\mu$s  & Full right limit \\
$45^\circ$  & $995\,\mu$s  & Right of centre \\
$90^\circ$  & $1490\,\mu$s & Centre (faces forward) \\
$135^\circ$ & $1985\,\mu$s & Left of centre \\
$180^\circ$ & $2480\,\mu$s & Full left limit \\
\bottomrule
\end{tabular}
\end{center}

\textbf{Why 30 iterations in the loop?} The \texttt{servo\_appointed\_detection}
function calls \texttt{servo\_pulse} 30 times, generating 30 PWM cycles.
Each cycle is $\approx 20$\,ms, so the total drive time is $30 \times 20 = 600$\,ms.
This ensures the servo physically reaches and settles at the target angle before
the function returns.

% ------ 2.2 Original functions ------
\subsection{Original Servo Functions (Existing Code)}

\begin{lstlisting}[style=arduino,caption={Original servo\_pulse and servo\_appointed\_detection}]
// Generates one PWM pulse of the correct width for the given angle
void servo_pulse(int ServoPin, int myangle)
{
  int PulseWidth;
  PulseWidth = (myangle * 11) + 500;  // Map angle (0-180) to pulse (500-2480 us)

  digitalWrite(ServoPin, HIGH);       // Start of HIGH pulse
  delayMicroseconds(PulseWidth);      // Hold HIGH for pulse width duration
  digitalWrite(ServoPin, LOW);        // End of HIGH pulse
  delay(20 - PulseWidth / 1000);      // Fill remainder of 20 ms period
}

// Moves servo to 'pos' degrees and holds it (30 PWM cycles = ~600 ms)
void servo_appointed_detection(int pos)
{
  int i = 0;
  for (i = 0; i <= 30; i++)
  {
    servo_pulse(ServoPin, pos);       // Repeat pulse to drive servo to position
  }
}
\end{lstlisting}

% ------ 2.3 Modified code ------
\subsection{Task 2 --- Modified Code: Specific Angle Test}

\begin{lstlisting}[style=arduino,caption={Task 2 modification --- ultrasonic servo at specific angles}]
// ================================================================
// TASK 2 -- Specific angle test routine for ultrasonic servo (J3)
// Define named positions matching the 5 test angles
// ================================================================

#define INITIAL_ANGLE   90   // Center: sensor faces directly forward
#define ANGLE_45R       45   // 45 degrees to the right of center
#define ANGLE_MAX_R      0   // Maximum right extent
#define ANGLE_45L      135   // 45 degrees to the left of center
#define ANGLE_MAX_L    180   // Maximum left extent
#define HOLD_TIME     2000   // Hold each position 2 seconds for measurement

void task2_ultrasonic_angles()
{
  // Test sequence: initial -> right -> full right -> left -> full left -> return
  int positions[] = {INITIAL_ANGLE, ANGLE_45R, ANGLE_MAX_R,
                     ANGLE_45L, ANGLE_MAX_L, INITIAL_ANGLE};
  const char* labels[] = {"INITIAL (90)", "45 Right", "MAX Right (0)",
                           "45 Left", "MAX Left (180)", "FINAL (90)"};
  int count = 6;

  whistle();   // Audible signal: test starting
  delay(500);

  for (int i = 0; i < count; i++)
  {
    // --- Move servo to target angle ---
    servo_appointed_detection(positions[i]);

    // --- LED color feedback by position type ---
    if (positions[i] == 90)                          // Center
      corlor_led(OFF, ON, OFF);   // Green
    else if (positions[i] == 0 || positions[i] == 180) // Limit
      corlor_led(ON, OFF, OFF);   // Red
    else                                             // Intermediate
      corlor_led(OFF, OFF, ON);   // Blue

    // --- Serial output for logging & verification ---
    Serial.print("[Task 2] ");
    Serial.print(labels[i]);
    Serial.print(" | Pulse width: ");
    Serial.print((positions[i] * 11) + 500);
    Serial.println(" us");

    // --- Hold position for physical measurement ---
    delay(HOLD_TIME);
  }

  whistle();   // Audible signal: test complete
  corlor_led(OFF, OFF, OFF);
  Serial.println("[Task 2] Ultrasonic servo test COMPLETE.");
}
\end{lstlisting}

\subsection{Task 2 --- Track Motor Position Demonstrations}

\begin{conceptbox}{DC Motor vs Servo: Key Difference}
DC motors (used for the tracks) do \emph{not} have position control.
They have \textbf{direction} and \textbf{speed} control via an H-bridge
circuit. Instead of angles, you program them with:
\begin{itemize}[leftmargin=1.5em]
  \item \textbf{Direction}: forward or backward (two digital pins per motor)
  \item \textbf{Speed}: PWM duty cycle, 0--255 via \texttt{analogWrite()}
\end{itemize}
The G1 Tank uses \emph{differential drive}: independent control of left and
right track speeds produces all navigation behaviours.
\end{conceptbox}

\begin{lstlisting}[style=arduino,caption={Task 2 --- DC motor position/state demonstration}]
// ================================================================
// TASK 2 -- Track motor state demonstration
// Shows all independent motion states (DOF 1 and DOF 2)
// ================================================================

void task2_track_motors()
{
  CarSpeedControl = 150;    // Medium speed (0-255 PWM range)
  Serial.println("[Task 2] Track motor demonstration starting...");

  // State 1: Both forward -- straight ahead
  Serial.println("State 1: FORWARD");
  run();   delay(1200);   brake();   delay(500);

  // State 2: Both backward -- straight back
  Serial.println("State 2: BACKWARD");
  back();  delay(1200);   brake();   delay(500);

  // State 3: Left backward, right forward -- spin left in place
  Serial.println("State 3: SPIN LEFT");
  spin_left();   delay(900);   brake();   delay(500);

  // State 4: Left forward, right backward -- spin right in place
  Serial.println("State 4: SPIN RIGHT");
  spin_right();  delay(900);   brake();   delay(500);

  // State 5: Left stop, right forward -- pivot around left track
  Serial.println("State 5: PIVOT LEFT");
  left();   delay(900);   brake();   delay(500);

  // State 6: Left forward, right stop -- pivot around right track
  Serial.println("State 6: PIVOT RIGHT");
  right();  delay(900);   brake();   delay(500);

  Serial.println("[Task 2] Track motor test COMPLETE.");
}
\end{lstlisting}

\subsection{Task 2 --- Camera Servo Code (Theoretical --- Pi Required)}

\begin{dangerbox}{BLOCKED --- Pi image failure}
The camera pan and tilt servos are wired to the Raspberry Pi GPIO header.
The following Python code \emph{would} execute these tests if the Pi were
operational. It is provided for completeness.
\end{dangerbox}

\begin{lstlisting}[style=python,caption={Task 2 camera servo code --- Pi version (theoretical)}]
import pigpio
import time

pi = pigpio.pi()          # Connect to local Pi GPIO daemon

PAN_PIN  = 2              # BCM GPIO pin for pan servo  (J4)
TILT_PIN = 3              # BCM GPIO pin for tilt servo (J5)

def angle_to_pw(angle):
    """Map 0-180 degrees to 500-2500 microsecond pulse width."""
    return int(500 + (angle / 180.0) * 2000)

def move_servo(pin, angle, label):
    pw = angle_to_pw(angle)
    pi.set_servo_pulsewidth(pin, pw)
    print(f"[Task 2] {label} -> {angle} deg | PW: {pw} us")
    time.sleep(2.0)          # Hold for physical measurement

def test_pan_servo():
    print("--- Pan servo (J4) test ---")
    for angle, label in [(90,"INITIAL"), (45,"45 right"), (0,"MAX right"),
                         (135,"45 left"), (180,"MAX left"), (90,"FINAL")]:
        move_servo(PAN_PIN, angle, f"Pan {label}")

def test_tilt_servo():
    print("--- Tilt servo (J5) test ---")
    for angle, label in [(90,"INITIAL horizontal"), (50,"Tilt up"),
                         (40,"MAX up"), (120,"Tilt down"),
                         (130,"MAX down"), (90,"FINAL horizontal")]:
        move_servo(TILT_PIN, angle, f"Tilt {label}")

test_pan_servo()
test_tilt_servo()
pi.stop()
\end{lstlisting}

% ============================================================
\newpage
\section{Task 3 --- Position Verification using Product of Exponentials (POE)}
% ============================================================

\begin{warnbox}{STATUS: PARTIAL --- J3 (ultrasonic) physical; J4, J5 (camera) theoretical}
The POE framework is applied to all five joints. Physical measurement against
POE predictions was possible only for joint $J_3$ (ultrasonic servo, UNO-accessible).
Camera joints $J_4$ and $J_5$ are verified analytically only.
\end{warnbox}

% ------ 3.1 What is POE ------
\subsection{What is the Product of Exponentials?}

\begin{conceptbox}{The Core Idea: Why POE?}
When a robot arm has multiple joints, you need to track where the
\emph{end-effector} (the tool at the tip) ends up after all joints move.
The classical approach (Denavit-Hartenberg) is error-prone to set up.
\textbf{Product of Exponentials (POE)} is the modern alternative:

\begin{itemize}[leftmargin=1.5em]
  \item Every joint rotation is expressed as a $4\times4$ matrix called a
        \textbf{matrix exponential}.
  \item You \emph{multiply} these matrices together (hence ``product'') along
        with the \textbf{home configuration} matrix $M$.
  \item The result is a $4\times4$ \textbf{transformation matrix} $T$ that
        tells you exactly where the end-effector is and how it is oriented.
\end{itemize}
\end{conceptbox}

\subsubsection{The Transformation Matrix T in SE(3)}

A $4 \times 4$ \textbf{homogeneous transformation matrix} encodes both
\emph{rotation} and \emph{translation} in one compact form:

\[
  T = \begin{pmatrix} \mat{R} & \vect{p} \\ \vect{0}^\top & 1 \end{pmatrix}
  \in \SE
\]

where $\mat{R} \in \SO$ is a $3\times3$ rotation matrix ($\mat{R}^\top\mat{R}=\mat{I}$,
$\det\mat{R}=+1$) and $\vect{p}\in\R^3$ is the position vector of the frame
origin.

\subsubsection{The Screw Axis $\mathcal{S}$}

\begin{conceptbox}{Screw Axis --- The Key Building Block}
Every revolute joint has a \textbf{screw axis} $\mathcal{S}$ that
encodes \emph{both} the axis of rotation and the geometry of where that
axis sits in space:
\[
  \mathcal{S} = \begin{pmatrix}\vect{\hat\omega} \\ \vect{v}\end{pmatrix}
  \in \R^6,
  \qquad
  \vect{v} = -\vect{\hat\omega} \times \vect{q}
\]

\begin{tabular}{cl}
  $\vect{\hat\omega}$ & Unit vector along the axis of rotation \\
  $\vect{q}$          & Any point on the axis of rotation (in the space frame) \\
  $\vect{v}$          & Linear velocity component (computed from $\vect{\hat\omega}$ and $\vect{q}$) \\
\end{tabular}

For a \emph{pure revolute} joint (no translation), the pitch is zero, so
$\mathcal{S}$ is determined entirely by $\vect{\hat\omega}$ and any point
$\vect{q}$ on the axis.
\end{conceptbox}

\subsubsection{The Matrix Exponential $e^{[\mathcal{S}]\theta}$}

The screw axis $\mathcal{S}$ is converted to a $6\times6$
\textbf{bracket} form $[\mathcal{S}]$, then exponentiated. For a revolute
joint (pure rotation), the $4\times4$ matrix exponential is computed via the
\textbf{Rodrigues rotation formula}:

\begin{mathbox}{Rodrigues Rotation Formula}
For unit axis $\vect{\hat\omega}$ and angle $\theta$:
\[
  \mat{R}(\vect{\hat\omega},\theta)
  = \mat{I}
  + \sin\theta\,[\vect{\hat\omega}]
  + (1 - \cos\theta)\,[\vect{\hat\omega}]^2
\]

where $[\vect{\hat\omega}]$ is the \textbf{skew-symmetric matrix} of
$\vect{\hat\omega} = (\omega_1,\omega_2,\omega_3)$:
\[
  [\vect{\hat\omega}]
  = \begin{pmatrix}
      0        & -\omega_3 &  \omega_2 \\
      \omega_3 &  0        & -\omega_1 \\
     -\omega_2 &  \omega_1 &  0
    \end{pmatrix}
\]

The full $4\times4$ matrix exponential is:
\[
  e^{[\mathcal{S}]\theta}
  = \begin{pmatrix}
      \mat{R} & (\mat{I}-\mat{R})\,[\vect{\hat\omega}]\vect{v}
              + \vect{\hat\omega}\,\vect{\hat\omega}^\top\vect{v}\,\theta \\
      \vect{0}^\top & 1
    \end{pmatrix}
\]

For a joint whose axis passes through the origin, the translation
simplifies to $\vect{p} = \vect{0}$.
\end{mathbox}

\subsection{The POE Forward Kinematics Formula}

\begin{mathbox}{Product of Exponentials --- Space Form}
\[
  T(\vect{\theta})
  = e^{[\mathcal{S}_1]\theta_1}
    \cdot e^{[\mathcal{S}_2]\theta_2}
    \cdots
    e^{[\mathcal{S}_n]\theta_n}
    \cdot M
\]

\begin{tabular}{cl}
$M$                    & Home configuration matrix (all $\theta_i = 0$) \\
$\mathcal{S}_i$        & Screw axis of joint $i$ in the \emph{space frame at home config} \\
$\theta_i$             & Angle of joint $i$ \\
$e^{[\mathcal{S}_i]\theta_i}$ & $4\times4$ matrix exponential \\
\end{tabular}

\textbf{Key insight:} Each exponential is computed using the joint's
geometry as it was \emph{at the home configuration} --- you do not need
to track how previous joints moved first. This is what makes POE elegant.
\end{mathbox}

\subsection{Space Frame Definition}

The \textbf{space frame} $\{S\}$ is fixed at the chassis center (ground link $L_1$):

\begin{center}
\renewcommand{\arraystretch}{1.2}
\begin{tabular}{ll}
\toprule
\textbf{Axis} & \textbf{Direction} \\
\midrule
$\hat{x}$ & Pointing forward (direction of travel) \\
$\hat{y}$ & Pointing left (lateral) \\
$\hat{z}$ & Pointing upward (vertical) \\
\midrule
\multicolumn{2}{l}{Origin at geometric center of chassis, at ground level} \\
\bottomrule
\end{tabular}
\end{center}

\vspace{6pt}
\textbf{Estimated dimensions from manual and physical measurement:}

\begin{center}
\begin{tabular}{lcc}
\toprule
\textbf{Dimension} & \textbf{Symbol} & \textbf{Value (m)} \\
\midrule
Ultrasonic servo axis forward offset & $q_{3x}$ & $0.08$ \\
Ultrasonic servo axis height & $q_{3z}$ & $0.08$ \\
Camera pan axis forward offset & $q_{4x}$ & $0.04$ \\
Camera pan axis height & $q_{4z}$ & $0.12$ \\
Camera tilt axis forward offset & $q_{5x}$ & $0.08$ \\
Camera tilt axis height & $q_{5z}$ & $0.15$ \\
Sensor tip from pan axis & $d_s$ & $0.045$ \\
\bottomrule
\end{tabular}
\end{center}

\subsection{Screw Axis Calculations}

\subsubsection{Joint $J_3$ --- Ultrasonic Servo}

\begin{mathbox}{Screw Axis $\mathcal{S}_3$}
Joint $J_3$ rotates about the vertical ($z$) axis:
\[
  \vect{\hat\omega}_3 = \begin{pmatrix}0\\0\\1\end{pmatrix}, \qquad
  \vect{q}_3 = \begin{pmatrix}0.08\\0\\0.08\end{pmatrix}
\]
\[
  \vect{v}_3
  = -\vect{\hat\omega}_3 \times \vect{q}_3
  = -\begin{pmatrix}0\\0\\1\end{pmatrix} \times \begin{pmatrix}0.08\\0\\0.08\end{pmatrix}
  = -\begin{pmatrix}(0)(0.08)-(1)(0)\\(1)(0.08)-(0)(0.08)\\(0)(0)-(0)(0.08)\end{pmatrix}
  = -\begin{pmatrix}0\\0.08\\0\end{pmatrix}
  = \begin{pmatrix}0\\-0.08\\0\end{pmatrix}
\]
\[
  \boxed{\mathcal{S}_3 = (0,\;0,\;1,\;0,\;-0.08,\;0)^\top}
\]
\end{mathbox}

\subsubsection{Joint $J_4$ --- Camera Pan Servo}

\begin{mathbox}{Screw Axis $\mathcal{S}_4$}
Joint $J_4$ also rotates about the vertical ($z$) axis:
\[
  \vect{\hat\omega}_4 = \begin{pmatrix}0\\0\\1\end{pmatrix}, \qquad
  \vect{q}_4 = \begin{pmatrix}0.04\\0\\0.12\end{pmatrix}
\]
\[
  \vect{v}_4
  = -\begin{pmatrix}0\\0\\1\end{pmatrix} \times \begin{pmatrix}0.04\\0\\0.12\end{pmatrix}
  = -\begin{pmatrix}0\\0.04\\0\end{pmatrix}
  = \begin{pmatrix}0\\-0.04\\0\end{pmatrix}
\]
\[
  \boxed{\mathcal{S}_4 = (0,\;0,\;1,\;0,\;-0.04,\;0)^\top}
\]
\end{mathbox}

\subsubsection{Joint $J_5$ --- Camera Tilt Servo}

\begin{mathbox}{Screw Axis $\mathcal{S}_5$}
Joint $J_5$ rotates about the lateral ($y$) axis (tilts in the $xz$-plane):
\[
  \vect{\hat\omega}_5 = \begin{pmatrix}0\\1\\0\end{pmatrix}, \qquad
  \vect{q}_5 = \begin{pmatrix}0.08\\0\\0.15\end{pmatrix}
\]
\[
  \vect{v}_5
  = -\begin{pmatrix}0\\1\\0\end{pmatrix} \times \begin{pmatrix}0.08\\0\\0.15\end{pmatrix}
  = -\begin{pmatrix}(1)(0.15)-(0)(0)\\(0)(0.08)-(0)(0.15)\\(0)(0)-(1)(0.08)\end{pmatrix}
  = -\begin{pmatrix}0.15\\0\\-0.08\end{pmatrix}
  = \begin{pmatrix}-0.15\\0\\0.08\end{pmatrix}
\]
\[
  \boxed{\mathcal{S}_5 = (0,\;1,\;0,\;-0.15,\;0,\;0.08)^\top}
\]
\end{mathbox}

\subsection{Skew-Symmetric Matrices for Each Joint Axis}

\begin{mathbox}{Bracket Forms $[\vect{\hat\omega}]$}
\textbf{For $J_3$ and $J_4$} (both rotate about $z$, $\vect{\hat\omega} = (0,0,1)$):
\[
  [\vect{\hat\omega}_{3,4}]
  = \begin{pmatrix}0&-1&0\\1&0&0\\0&0&0\end{pmatrix},
  \qquad
  [\vect{\hat\omega}_{3,4}]^2
  = \begin{pmatrix}-1&0&0\\0&-1&0\\0&0&0\end{pmatrix}
\]

\textbf{For $J_5$} (rotates about $y$, $\vect{\hat\omega} = (0,1,0)$):
\[
  [\vect{\hat\omega}_5]
  = \begin{pmatrix}0&0&1\\0&0&0\\-1&0&0\end{pmatrix},
  \qquad
  [\vect{\hat\omega}_5]^2
  = \begin{pmatrix}-1&0&0\\0&0&0\\0&0&-1\end{pmatrix}
\]
\end{mathbox}

\subsection{POE Calculation Examples}

\subsubsection{Ultrasonic Sensor at $\theta_3 = 90^\circ$ (Faces Left)}

\begin{mathbox}{Full calculation: $e^{[\mathcal{S}_3]\,90^\circ}$}
\textbf{Step 1: Rotation matrix} $\mat{R}_3(90^\circ)$ via Rodrigues:
\begin{align*}
  \mat{R}_3
  &= \mat{I} + \sin(90^\circ)[\vect{\hat\omega}_3] + (1-\cos(90^\circ))[\vect{\hat\omega}_3]^2 \\
  &= \mat{I} + 1.0\begin{pmatrix}0&-1&0\\1&0&0\\0&0&0\end{pmatrix}
           + 1.0\begin{pmatrix}-1&0&0\\0&-1&0\\0&0&0\end{pmatrix} \\[6pt]
  &= \begin{pmatrix}1&0&0\\0&1&0\\0&0&1\end{pmatrix}
   + \begin{pmatrix}0&-1&0\\1&0&0\\0&0&0\end{pmatrix}
   + \begin{pmatrix}-1&0&0\\0&-1&0\\0&0&0\end{pmatrix}
   = \begin{pmatrix}0&-1&0\\1&0&0\\0&0&1\end{pmatrix}
\end{align*}

\textbf{Step 2: Translation vector} $\vect{p}_3$:
\[
  \vect{p}_3
  = (\mat{I}-\mat{R}_3)[\vect{\hat\omega}_3]\vect{v}_3
  = \begin{pmatrix}1&1&0\\-1&1&0\\0&0&0\end{pmatrix}
    \begin{pmatrix}0&-1&0\\1&0&0\\0&0&0\end{pmatrix}
    \begin{pmatrix}0\\-0.08\\0\end{pmatrix}
  = \begin{pmatrix}0.08\\-0.08\\0\end{pmatrix}
\]

\textbf{Step 3: Full transformation matrix:}
\[
  e^{[\mathcal{S}_3]\,90^\circ}
  = \begin{pmatrix}0&-1&0&0.08\\1&0&0&-0.08\\0&0&1&0\\0&0&0&1\end{pmatrix}
\]

\textbf{Interpretation:} The sensor has rotated $90^\circ$ CCW (now faces left),
and its mounting point has shifted $+8$\,cm in $x$ and $-8$\,cm in $y$
from the chassis center.
\end{mathbox}

\subsubsection{Ultrasonic Sensor at $\theta_3 = 45^\circ$ (Faces $45^\circ$ Right)}

\begin{mathbox}{Calculation: $e^{[\mathcal{S}_3]\,45^\circ}$}
\[
  \sin(45^\circ) = \cos(45^\circ) = \frac{\sqrt{2}}{2} \approx 0.7071
\]

\textbf{Rotation matrix:}
\[
  \mat{R}_3(45^\circ)
  = \mat{I} + 0.7071[\vect{\hat\omega}_3] + (1-0.7071)[\vect{\hat\omega}_3]^2
  = \begin{pmatrix}0.7071&-0.7071&0\\0.7071&0.7071&0\\0&0&1\end{pmatrix}
\]

\textbf{Translation:}
\[
  p_x = (1-\cos 45^\circ)\times 0.08 = 0.2929 \times 0.08 = 0.0234\text{ m}
\]
\[
  p_y = \sin(45^\circ) \times 0.08 = 0.7071 \times 0.08 = 0.0566\text{ m}
\]

\textbf{Full matrix:}
\[
  e^{[\mathcal{S}_3]\,45^\circ}
  = \begin{pmatrix}
      0.707 & -0.707 & 0 & 0.023 \\
      0.707 &  0.707 & 0 & 0.057 \\
      0     &  0     & 1 & 0     \\
      0     &  0     & 0 & 1
    \end{pmatrix}
\]

\textbf{Physical check:} Sensor tip moves
$d\sin(45^\circ) = 0.08 \times 0.707 \approx 5.7$\,cm to the left.
This was measurable with a ruler during the lab.
\end{mathbox}

\subsubsection{Camera Tilt at $\theta_5 = 40^\circ$ Upward (Theoretical)}

\begin{mathbox}{Calculation: $e^{[\mathcal{S}_5]\,40^\circ}$}
\[
  \sin(40^\circ) = 0.6428, \quad (1-\cos 40^\circ) = 0.2340
\]

\textbf{Rotation matrix:}
\[
  \mat{R}_5(40^\circ)
  = \mat{I}
  + 0.6428\begin{pmatrix}0&0&1\\0&0&0\\-1&0&0\end{pmatrix}
  + 0.2340\begin{pmatrix}-1&0&0\\0&0&0\\0&0&-1\end{pmatrix}
  = \begin{pmatrix}0.766&0&0.643\\0&1&0\\-0.643&0&0.766\end{pmatrix}
\]

\textbf{Interpretation:} Camera has tilted upward by $40^\circ$ in the $xz$-plane.
The forward component of the camera direction decreases while an upward
component emerges.
\end{mathbox}

\subsection{Home Configuration Matrices $M$}

\begin{mathbox}{Home configuration for the ultrasonic sensor (all joints = 0)}
At the home position, the ultrasonic sensor faces forward. Estimated geometry:
\[
  M_{\text{ultrasonic}}
  = \begin{pmatrix}
      1 & 0 & 0 & 0.125 \\
      0 & 1 & 0 & 0.000 \\
      0 & 0 & 1 & 0.080 \\
      0 & 0 & 0 & 1
    \end{pmatrix}
\]
where $p_x = 0.08 + 0.045 = 0.125$\,m (servo offset + sensor body),
$p_z = 0.08$\,m (servo height above chassis).
\end{mathbox}

\subsection{Physical Verification Procedure}

\begin{donebox}{How to Physically Verify POE Results (Without Camera)}
\begin{enumerate}[leftmargin=1.5em]
  \item Place robot on a flat surface. Mark the chassis center and forward
        direction on a sheet of paper below.
  \item Tape a protractor to the chassis plate, aligned with the servo
        rotation axis.
  \item Command each angle via \texttt{servo\_appointed\_detection(angle)};
        trigger from the serial monitor or a test loop.
  \item Read the physical angle from the protractor. Compare to the
        commanded angle. Tolerance: $\pm 3^\circ$ for software PWM.
  \item For translation: measure the sensor tip position with a ruler
        from the joint axis, and compare to the POE-predicted position.
  \item Record \texttt{Serial.print()} output as the digital log.
\end{enumerate}
\end{donebox}

\begin{table}[h!]
\centering
\caption{Commanded vs predicted positions for $J_3$ (ultrasonic servo)}
\renewcommand{\arraystretch}{1.3}
\begin{tabularx}{\linewidth}{C{1.8cm} C{2cm} C{2cm} L{4cm} C{2.8cm}}
\toprule
\rowcolor{MidBlue}
\color{white}\textbf{Angle} &
\color{white}\textbf{PW ($\mu$s)} &
\color{white}\textbf{$p_y$ shift} &
\color{white}\textbf{POE Prediction} &
\color{white}\textbf{Verified} \\
\midrule
$0^\circ$   & 500  & $0$\,m   & Faces $90^\circ$ right & \cellcolor{SuccessBG}\color{SuccessGreen}Yes \\
$45^\circ$  & 995  & $+0.057$\,m & $45^\circ$ right, tip $5.7$\,cm left & \cellcolor{SuccessBG}\color{SuccessGreen}Yes \\
$90^\circ$  & 1490 & $+0.080$\,m & Faces forward, tip at origin & \cellcolor{SuccessBG}\color{SuccessGreen}Yes \\
$135^\circ$ & 1985 & $+0.057$\,m & $45^\circ$ left & \cellcolor{SuccessBG}\color{SuccessGreen}Yes \\
$180^\circ$ & 2480 & $0$\,m   & Faces $90^\circ$ left & \cellcolor{SuccessBG}\color{SuccessGreen}Yes \\
\bottomrule
\end{tabularx}
\end{table}

% ============================================================
\newpage
\section{Task 4 --- Per-Actuator Position Programming \& Verification}
% ============================================================

% \begin{conceptbox}{What Task 4 Requires --- Background for Learners}
% Task 4 asks you to, for \emph{each individual actuator}:
% \begin{enumerate}[leftmargin=1.5em]
%   \item Define and program an \textbf{initial position} (the starting/home state).
%   \item Program \textbf{multiple intermediate positions} (showing the range of motion).
%   \item Define and program the \textbf{final position} (often returns to initial).
%   \item \textbf{Verify} each position --- physically (with protractor/ruler) where
%         possible, and using the POE formula as the mathematical ground truth.
% \end{enumerate}
% The goal is to demonstrate that you understand each actuator's range of motion
% and can control it precisely.
% \end{conceptbox}

% ------ Actuator 1 ------
\subsection{Actuator 1 --- Left DC Motor (DOF 1)}

\begin{donebox}{}
Left motor is fully accessible from Arduino UNO via pins 7, 8, and PWM pin 6.
\end{donebox}

\begin{conceptbox}{Understanding H-Bridge Motor Control}
A DC motor is driven through an \textbf{H-bridge} circuit (L298N on the
Yahboom expansion board). The H-bridge allows current to flow through the motor
in either direction by switching four transistors. The UNO controls it with:

\begin{itemize}[leftmargin=1.5em]
  \item \textbf{Direction pin A} (pin 8): HIGH = forward polarity
  \item \textbf{Direction pin B} (pin 7): HIGH = reverse polarity
  \item \textbf{PWM pin} (pin 6): \texttt{analogWrite(0--255)} sets speed
  \item Both direction pins LOW = brake (motor terminals shorted together)
\end{itemize}

Speed in RPM is proportional to PWM duty cycle. At \texttt{analogWrite(150)},
the motor receives approximately $150/255 = 59\%$ of battery voltage.
\end{conceptbox}

\vspace{6pt}
\begin{table}[h!]
\centering
\caption{Left DC motor states}
\renewcommand{\arraystretch}{1.3}
\begin{tabularx}{\linewidth}{L{2.5cm} C{2.5cm} C{2.5cm} C{2.5cm} L{3cm}}
\toprule
\rowcolor{MidBlue}
\color{white}\textbf{State} &
\color{white}\textbf{Pin 8 (go)} &
\color{white}\textbf{Pin 7 (back)} &
\color{white}\textbf{PWM (Pin 6)} &
\color{white}\textbf{Effect} \\
\midrule
Forward & HIGH & LOW & 0--255 & Track drives forward \\
Backward & LOW & HIGH & 0--255 & Track drives backward \\
Brake & LOW & LOW & 0 & Motor locked (short) \\
\bottomrule
\end{tabularx}
\end{table}

\begin{lstlisting}[style=arduino,caption={Task 4 --- Left motor position/state test}]
// Task 4: Left motor -- initial, multiple positions, final
void task4_left_motor()
{
  Serial.println("=== Task 4: Left Motor ===");

  // INITIAL: Stopped
  Serial.println("[INITIAL] Stopped");
  digitalWrite(Left_motor_go,   LOW);
  digitalWrite(Left_motor_back, LOW);
  analogWrite(Left_motor_pwm, 0);
  delay(1000);

  // POSITION 1: Slow forward (PWM = 80, ~31% duty cycle)
  Serial.println("[POS 1] Slow forward - PWM 80");
  CarSpeedControl = 80;
  digitalWrite(Left_motor_go,   HIGH);
  digitalWrite(Left_motor_back, LOW);
  analogWrite(Left_motor_pwm, CarSpeedControl);
  delay(1200);

  // POSITION 2: Medium forward (PWM = 150, ~59% duty cycle)
  Serial.println("[POS 2] Medium forward - PWM 150");
  CarSpeedControl = 150;
  analogWrite(Left_motor_pwm, CarSpeedControl);
  delay(1200);

  // POSITION 3: Fast forward (PWM = 220, ~86% duty cycle)
  Serial.println("[POS 3] Fast forward - PWM 220");
  CarSpeedControl = 220;
  analogWrite(Left_motor_pwm, CarSpeedControl);
  delay(1200);

  // POSITION 4: Medium backward
  Serial.println("[POS 4] Medium backward - PWM 150");
  CarSpeedControl = 150;
  digitalWrite(Left_motor_go,   LOW);
  digitalWrite(Left_motor_back, HIGH);
  analogWrite(Left_motor_pwm, CarSpeedControl);
  delay(1200);

  // FINAL: Stopped (return to initial)
  Serial.println("[FINAL] Stopped");
  digitalWrite(Left_motor_go,   LOW);
  digitalWrite(Left_motor_back, LOW);
  analogWrite(Left_motor_pwm, 0);
  delay(500);
  Serial.println("=== Left Motor Test COMPLETE ===");
}
\end{lstlisting}

% ------ Actuator 2 ------
\subsection{Actuator 2 --- Right DC Motor (DOF 2)}

\begin{donebox}{}
Mirror of left motor. Pins: go=2, back=4, PWM=5.
\end{donebox}

\begin{lstlisting}[style=arduino,caption={Task 4 --- Right motor position/state test}]
// Task 4: Right motor -- initial, multiple positions, final
void task4_right_motor()
{
  Serial.println("=== Task 4: Right Motor ===");

  // INITIAL: Stopped
  Serial.println("[INITIAL] Stopped");
  digitalWrite(Right_motor_go,   LOW);
  digitalWrite(Right_motor_back, LOW);
  analogWrite(Right_motor_pwm, 0);
  delay(1000);

  // POSITION 1: Slow forward
  Serial.println("[POS 1] Slow forward - PWM 80");
  CarSpeedControl = 80;
  digitalWrite(Right_motor_go,   HIGH);
  digitalWrite(Right_motor_back, LOW);
  analogWrite(Right_motor_pwm, CarSpeedControl);
  delay(1200);

  // POSITION 2: Medium forward
  Serial.println("[POS 2] Medium forward - PWM 150");
  CarSpeedControl = 150;
  analogWrite(Right_motor_pwm, CarSpeedControl);
  delay(1200);

  // POSITION 3: Fast forward
  Serial.println("[POS 3] Fast forward - PWM 220");
  CarSpeedControl = 220;
  analogWrite(Right_motor_pwm, CarSpeedControl);
  delay(1200);

  // POSITION 4: Medium backward
  Serial.println("[POS 4] Medium backward - PWM 150");
  CarSpeedControl = 150;
  digitalWrite(Right_motor_go,   LOW);
  digitalWrite(Right_motor_back, HIGH);
  analogWrite(Right_motor_pwm, CarSpeedControl);
  delay(1200);

  // FINAL: Stopped
  Serial.println("[FINAL] Stopped");
  digitalWrite(Right_motor_go,   LOW);
  digitalWrite(Right_motor_back, LOW);
  analogWrite(Right_motor_pwm, 0);
  delay(500);
  Serial.println("=== Right Motor Test COMPLETE ===");
}
\end{lstlisting}

% ------ Actuator 3 ------
\subsection{Actuator 3 --- Ultrasonic Servo, Joint $J_3$ (DOF 3)}

\begin{donebox}{}
All five test angles plus return-to-home were physically executed and verified
with a protractor. POE predictions match observed positions within $\pm 3^\circ$.
\end{donebox}

\begin{table}[h!]
\centering
\caption{Ultrasonic servo --- all programmed positions}
\renewcommand{\arraystretch}{1.3}
\begin{tabularx}{\linewidth}{L{2.5cm} C{1.5cm} C{1.8cm} C{2cm} L{3.5cm} C{2cm}}
\toprule
\rowcolor{MidBlue}
\color{white}\textbf{Position} &
\color{white}\textbf{Angle} &
\color{white}\textbf{PW ($\mu$s)} &
\color{white}\textbf{LED Color} &
\color{white}\textbf{Physical State} &
\color{white}\textbf{Verified} \\
\midrule
Initial (Home) & $90^\circ$ & 1490 & Green & Sensor faces forward & \cellcolor{SuccessBG}\color{SuccessGreen}Yes \\
Position 1 & $45^\circ$ & 995 & Blue & $45^\circ$ right of fwd & \cellcolor{SuccessBG}\color{SuccessGreen}Yes \\
Position 2 (Min) & $0^\circ$ & 500 & Red & Full right limit & \cellcolor{SuccessBG}\color{SuccessGreen}Yes \\
Position 3 & $135^\circ$ & 1985 & Blue & $45^\circ$ left of fwd & \cellcolor{SuccessBG}\color{SuccessGreen}Yes \\
Position 4 (Max) & $180^\circ$ & 2480 & Red & Full left limit & \cellcolor{SuccessBG}\color{SuccessGreen}Yes \\
Final (Return) & $90^\circ$ & 1490 & Green & Returns to forward & \cellcolor{SuccessBG}\color{SuccessGreen}Yes \\
\bottomrule
\end{tabularx}
\end{table}

\begin{lstlisting}[style=arduino,caption={Task 4 --- Ultrasonic servo: initial, 4 positions, final with logging}]
void task4_ultrasonic_servo()
{
  Serial.println("=== Task 4: Ultrasonic Servo (J3) ===");
  int angles[] = { 90,  45,   0, 135, 180,  90 };
  int pw[]     = {1490, 995, 500,1985,2480, 1490};
  const char* names[] = {"INITIAL (90)", "POS1 (45)",
                          "POS2-MIN (0)", "POS3 (135)",
                          "POS4-MAX (180)", "FINAL (90)"};
  int n = 6;

  whistle();
  delay(400);

  for (int i = 0; i < n; i++)
  {
    servo_appointed_detection(angles[i]);

    // LED color: green=center, red=limit, blue=intermediate
    if      (angles[i] == 90)                corlor_led(OFF, ON,  OFF);
    else if (angles[i] == 0 || angles[i]==180) corlor_led(ON,  OFF, OFF);
    else                                      corlor_led(OFF, OFF, ON);

    Serial.print("["); Serial.print(names[i]); Serial.print("] ");
    Serial.print("Commanded: "); Serial.print(angles[i]); Serial.print(" deg | ");
    Serial.print("PW: "); Serial.print(pw[i]); Serial.println(" us");

    delay(2500);   // Hold 2.5 s for protractor reading
  }

  whistle();
  corlor_led(OFF, OFF, OFF);
  Serial.println("=== Ultrasonic Servo Test COMPLETE ===");
}
\end{lstlisting}

% ------ Actuator 4 ------
\subsection{Actuator 4 --- Camera Pan Servo, Joint $J_4$ (DOF 4)}

\begin{dangerbox}{STATUS: THEORETICAL ONLY --- Pi GPIO inaccessible from UNO}
Camera pan servo is wired to the Raspberry Pi GPIO, not the Arduino UNO.
Physical testing was not possible due to the Pi image failure.
Full mathematical analysis is provided below.
\end{dangerbox}

\begin{table}[h!]
\centering
\caption{Camera pan servo --- theoretical position table}
\renewcommand{\arraystretch}{1.3}
\begin{tabularx}{\linewidth}{L{2.5cm} C{1.5cm} C{1.8cm} L{3.5cm} C{2.8cm}}
\toprule
\rowcolor{MidBlue}
\color{white}\textbf{Position} &
\color{white}\textbf{Angle} &
\color{white}\textbf{PW ($\mu$s)} &
\color{white}\textbf{Physical State} &
\color{white}\textbf{Status} \\
\midrule
Initial (Home) & $90^\circ$ & 1490 & Camera faces forward & \cellcolor{DangerBG}\color{DangerRed}Theoretical \\
Position 1 & $45^\circ$ & 995 & $45^\circ$ right & \cellcolor{DangerBG}\color{DangerRed}Theoretical \\
Position 2 (Min) & $0^\circ$ & 500 & Full right limit & \cellcolor{DangerBG}\color{DangerRed}Theoretical \\
Position 3 & $135^\circ$ & 1985 & $45^\circ$ left & \cellcolor{DangerBG}\color{DangerRed}Theoretical \\
Position 4 (Max) & $180^\circ$ & 2480 & Full left limit & \cellcolor{DangerBG}\color{DangerRed}Theoretical \\
Final (Return) & $90^\circ$ & 1490 & Returns forward & \cellcolor{DangerBG}\color{DangerRed}Theoretical \\
\bottomrule
\end{tabularx}
\end{table}

\textbf{POE at $\theta_4 = 90^\circ$:}
Because $J_4$ has the same axis direction as $J_3$ (both rotate about $z$), the
rotation matrix is identical:
\[
  \mat{R}_4(90^\circ) = \begin{pmatrix}0&-1&0\\1&0&0\\0&0&1\end{pmatrix}
\]
The translation differs because $\vect{q}_4 = (0.04, 0, 0.12)$:
\[
  \vect{p}_4 = (\mat{I}-\mat{R}_4)[\vect{\hat\omega}_4]\vect{v}_4
  = \begin{pmatrix}0.04\\-0.04\\0\end{pmatrix}
\]
\[
  e^{[\mathcal{S}_4]\,90^\circ}
  = \begin{pmatrix}0&-1&0&0.04\\1&0&0&-0.04\\0&0&1&0\\0&0&0&1\end{pmatrix}
\]

% ------ Actuator 5 ------
\subsection{Actuator 5 --- Camera Tilt Servo, Joint $J_5$ (DOF 5)}

\begin{dangerbox}{STATUS: THEORETICAL ONLY --- Pi GPIO inaccessible from UNO}
As with the pan servo, physical testing was not possible. Full analytical
treatment is provided.
\end{dangerbox}

\begin{conceptbox}{About the Tilt Servo Range}
Unlike the pan servo (which is mechanically unconstrained in software from
$0^\circ$ to $180^\circ$), the tilt servo is \emph{mechanically limited} by the
camera bracket geometry. From the manual drawings and physical inspection:
\begin{itemize}[leftmargin=1.5em]
  \item The bracket prevents tilting more than $\approx 40^\circ$ upward
        from horizontal (software $\approx 50^\circ$).
  \item It prevents tilting more than $\approx 30^\circ$ downward
        from horizontal (software $\approx 120^\circ$).
  \item Commanding beyond these limits will stall the servo and may damage it.
\end{itemize}
Always use the serial monitor to increment in small steps when testing a new
servo range.
\end{conceptbox}

\begin{table}[h!]
\centering
\caption{Camera tilt servo --- theoretical position table}
\renewcommand{\arraystretch}{1.3}
\begin{tabularx}{\linewidth}{L{2.5cm} C{1.5cm} C{1.8cm} L{3.5cm} C{2.8cm}}
\toprule
\rowcolor{MidBlue}
\color{white}\textbf{Position} &
\color{white}\textbf{Angle} &
\color{white}\textbf{PW ($\mu$s)} &
\color{white}\textbf{Physical State} &
\color{white}\textbf{Status} \\
\midrule
Initial (Home) & $90^\circ$ & 1490 & Camera horizontal, forward-facing & \cellcolor{DangerBG}\color{DangerRed}Theoretical \\
Tilt up $40^\circ$ & $50^\circ$ & 1050 & Camera looks $40^\circ$ upward & \cellcolor{DangerBG}\color{DangerRed}Theoretical \\
Max up (limit) & $\approx40^\circ$ & $\approx940$ & Mechanical upper limit & \cellcolor{DangerBG}\color{DangerRed}Theoretical \\
Tilt down $30^\circ$ & $120^\circ$ & 1820 & Camera looks $30^\circ$ downward & \cellcolor{DangerBG}\color{DangerRed}Theoretical \\
Max down (limit) & $\approx130^\circ$ & $\approx1930$ & Mechanical lower limit & \cellcolor{DangerBG}\color{DangerRed}Theoretical \\
Final (Return) & $90^\circ$ & 1490 & Returns horizontal & \cellcolor{DangerBG}\color{DangerRed}Theoretical \\
\bottomrule
\end{tabularx}
\end{table}

\textbf{POE verification at $\theta_5 = 40^\circ$ upward:}
From Section~3.6:
\[
  e^{[\mathcal{S}_5]\,40^\circ}
  = \begin{pmatrix}0.766&0&0.643&p_x\\0&1&0&0\\-0.643&0&0.766&p_z\\0&0&0&1\end{pmatrix}
\]

The translation $\vect{p}$ accounts for the offset of $J_5$ from the space frame origin.

% ------ 4.7 Summary ------
\subsection{Summary Table --- All Actuators}

\begin{table}[h!]
\centering
\caption{Complete per-actuator position summary}
\renewcommand{\arraystretch}{1.4}
\begin{tabularx}{\linewidth}{L{3cm} C{1.5cm} C{1.8cm} C{1.8cm} C{1.8cm} C{1.8cm} C{2cm}}
\toprule
\rowcolor{MidBlue}
\color{white}\textbf{Actuator} &
\color{white}\textbf{Type} &
\color{white}\textbf{Initial} &
\color{white}\textbf{Min} &
\color{white}\textbf{Max} &
\color{white}\textbf{Final} &
\color{white}\textbf{Status} \\
\midrule
Left DC Motor & DC Motor & Stop & PWM=0 & PWM=255 & Stop & \cellcolor{SuccessBG}\color{SuccessGreen}Physical \\
Right DC Motor & DC Motor & Stop & PWM=0 & PWM=255 & Stop & \cellcolor{SuccessBG}\color{SuccessGreen}Physical \\
Ultrasonic Servo ($J_3$) & Servo & $90^\circ$ & $0^\circ$ & $180^\circ$ & $90^\circ$ & \cellcolor{SuccessBG}\color{SuccessGreen}Physical \\
Camera Pan ($J_4$) & Servo & $90^\circ$ & $0^\circ$ & $180^\circ$ & $90^\circ$ & \cellcolor{DangerBG}\color{DangerRed}Theoretical \\
Camera Tilt ($J_5$) & Servo & $90^\circ$ & ${\approx}40^\circ$ & ${\approx}130^\circ$ & $90^\circ$ & \cellcolor{DangerBG}\color{DangerRed}Theoretical \\
\bottomrule
\end{tabularx}
\end{table}

% ============================================================
\newpage
\section{Conclusion}
% ============================================================


\vspace{8pt}
The Raspberry Pi OS image failure meant that camera pan and tilt servos (DOF 4
and DOF 5, joints $J_4$ and $J_5$) could not be physically operated or verified.
All required mathematical analysis --- screw axis derivation, Rodrigues rotation
matrices, matrix exponentials, and POE forward kinematics --- was completed in
full for these joints using the robot geometry from the manual. DOF 1, 2, and 3
(left track, right track, and ultrasonic servo) were fully demonstrated physically
and verified against POE predictions.

% ============================================================
\newpage
\appendix
\section{POE Deep Dive --- Chain Composition and Full Robot Analysis}
% ============================================================

\begin{conceptbox}{Why This Section Exists}
The earlier POE sections treated each joint in isolation. In a real robot, the
joints are \emph{chained} --- the motion of $J_4$ feeds directly into where
$J_5$ sits. This section explains how the full chain works, works through a
complete multi-joint product, and shows what each part of the resulting
transformation matrix physically means.
\end{conceptbox}

\subsection{Understanding Matrix Multiplication Order in POE}

A critical point that confuses beginners: in the POE \emph{space form}, all
screw axes $\mathcal{S}_i$ are defined in the \textbf{fixed space frame at the
home configuration} --- not in the current (moved) frame. This means:

\begin{resultbox}
\centering
You compute \emph{every} screw axis \textbf{once}, before any joint moves.
You never have to update or re-express them as joints rotate.
This is the main advantage of POE over the Denavit-Hartenberg method.
\end{resultbox}

\vspace{6pt}

The multiplication order in:
\[
  T = e^{[\mathcal{S}_1]\theta_1}
    \cdot e^{[\mathcal{S}_2]\theta_2}
    \cdot e^{[\mathcal{S}_3]\theta_3}
    \cdot e^{[\mathcal{S}_4]\theta_4}
    \cdot e^{[\mathcal{S}_5]\theta_5}
    \cdot M
\]
means: apply joint 1 first (leftmost), then joint 2, \ldots, then joint 5, then
apply the home offset $M$. Reading right to left gives the physical sequence
from end-effector back to base.

\subsection{Kinematic Chain Structure}

\textbf{Important observation:} $J_1$ and $J_2$ are \emph{parallel independent
branches} from the chassis. $J_4$ and $J_5$ are \emph{serial}: $J_5$ is mounted
on the output of $J_4$, so both must be in the chain to find the camera position.

\begin{center}
\renewcommand{\arraystretch}{1.2}
\begin{tabular}{L{3cm} L{11cm}}
\toprule
\textbf{Chain} & \textbf{Joints involved} \\
\midrule
Left track      & $J_1$ only (branch from chassis) \\
Right track     & $J_2$ only (branch from chassis) \\
Ultrasonic sensor & $J_3$ only (branch from chassis) \\
Camera          & $J_4 \to J_5$ in series (chassis $\to$ pan bracket $\to$ camera) \\
\bottomrule
\end{tabular}
\end{center}

\subsection{Full Camera Chain --- Joints $J_4$ and $J_5$ in Series}

For the camera end-effector at $L_6$:
\[
  T_{\text{camera}}(\theta_4,\,\theta_5)
  = e^{[\mathcal{S}_4]\theta_4}
    \cdot e^{[\mathcal{S}_5]\theta_5}
    \cdot M_{\text{camera}}
\]

\begin{mathbox}{Worked Example: Camera at $\theta_4 = 45^\circ$, $\theta_5 = 0^\circ$}

\textbf{Step 1:} Compute $e^{[\mathcal{S}_4]\,45^\circ}$.
$\mathcal{S}_4$ uses $z$-axis with $\vect{q}_4=(0.04,\,0,\,0.12)$:
\[
  \mat{R}_4(45^\circ) =
  \begin{pmatrix}0.707&-0.707&0\\0.707&0.707&0\\0&0&1\end{pmatrix},
  \quad
  \vect{p}_4 =
  \begin{pmatrix}(1-\cos45^\circ)(0.04)\\\sin45^\circ(0.04)\\0\end{pmatrix}
  =\begin{pmatrix}0.012\\0.028\\0\end{pmatrix}
\]
\[
  e^{[\mathcal{S}_4]\,45^\circ}
  = \begin{pmatrix}0.707&-0.707&0&0.012\\0.707&0.707&0&0.028\\0&0&1&0\\0&0&0&1\end{pmatrix}
\]

\textbf{Step 2:} $\theta_5=0^\circ$ so $e^{[\mathcal{S}_5]\,0^\circ}=\mathbf{I}_{4\times4}$.

\textbf{Step 3:} Home configuration for camera
($p_x=0.08$\,m, $p_z=0.15$\,m):
\[
  M_{\text{camera}}
  = \begin{pmatrix}1&0&0&0.08\\0&1&0&0\\0&0&1&0.15\\0&0&0&1\end{pmatrix}
\]

\textbf{Step 4:} Multiply chain left to right:
\begin{align*}
  T_{\text{camera}}
  &= \begin{pmatrix}0.707&-0.707&0&0.012\\0.707&0.707&0&0.028\\0&0&1&0\\0&0&0&1\end{pmatrix}
     \begin{pmatrix}1&0&0&0.08\\0&1&0&0\\0&0&1&0.15\\0&0&0&1\end{pmatrix} \\[8pt]
  &= \begin{pmatrix}
       0.707 & -0.707 & 0 & (0.707)(0.08)+(-0.707)(0)+0.012 \\
       0.707 &  0.707 & 0 & (0.707)(0.08)+(0.707)(0)+0.028  \\
       0     &  0     & 1 & 0.15 \\
       0     &  0     & 0 & 1
     \end{pmatrix} \\[8pt]
  &= \begin{pmatrix}
       0.707 & -0.707 & 0 & 0.069 \\
       0.707 &  0.707 & 0 & 0.085 \\
       0     &  0     & 1 & 0.150 \\
       0     &  0     & 0 & 1
     \end{pmatrix}
\end{align*}

\textbf{Physical meaning:} Camera has panned $45^\circ$ right. Optical centre
is at $(6.9\,\text{cm},\;8.5\,\text{cm},\;15\,\text{cm})$ relative to the
chassis centre. The rotation block shows the camera forward axis pointing
$45^\circ$ between the original forward and left directions.
\end{mathbox}

\subsection{Reading a Transformation Matrix}

\begin{mathbox}{Anatomy of $T$}
\[
  T
  = \begin{pmatrix}
      R_{11} & R_{12} & R_{13} & p_x \\
      R_{21} & R_{22} & R_{23} & p_y \\
      R_{31} & R_{32} & R_{33} & p_z \\
      0      & 0      & 0      & 1
    \end{pmatrix}
\]

\begin{center}
\renewcommand{\arraystretch}{1.3}
\begin{tabular}{L{3.5cm} L{9cm}}
\toprule
\textbf{Part} & \textbf{Meaning} \\
\midrule
Top-left $3\times3$ ($\mat{R}$) & Each column is a unit vector showing
  where the body's $x$, $y$, $z$ axes now point in the space frame \\
Column 1 of $\mat{R}$ & Direction the body's $x$-axis (forward) points \\
Column 2 of $\mat{R}$ & Direction the body's $y$-axis (left) points \\
Column 3 of $\mat{R}$ & Direction the body's $z$-axis (up) points \\
$(p_x,\,p_y,\,p_z)^\top$ & Position of the frame origin in the space frame \\
Bottom row $(0,0,0,1)$ & Homogeneous row --- always this, never changes \\
\midrule
\multicolumn{2}{l}{\textbf{Quick sanity check:} each column of $\mat{R}$ must have magnitude 1,}\\
\multicolumn{2}{l}{and any two columns must be perpendicular. Failure = arithmetic error.}\\
\bottomrule
\end{tabular}
\end{center}
\end{mathbox}

\subsection{POE Verification Checklist}

\begin{donebox}{Step-by-Step Checklist for Every POE Calculation}
\begin{enumerate}[leftmargin=1.5em,itemsep=2pt]
  \item \textbf{Identify} the end-effector and trace the kinematic chain from base to tip.
  \item \textbf{Define} the space frame origin and axes clearly before any maths.
  \item \textbf{Find} $\vect{\hat\omega}_i$ (unit axis vector) and $\vect{q}_i$
        (any point on the axis) for each joint in the chain.
  \item \textbf{Compute} $\vect{v}_i = -\vect{\hat\omega}_i \times \vect{q}_i$.
  \item \textbf{Form} the skew matrix $[\vect{\hat\omega}_i]$.
        Check: antisymmetric, zeros on diagonal.
  \item \textbf{Apply Rodrigues} to get $\mat{R}_i(\theta_i)$.
        Check: $\mat{R}^\top\mat{R}=\mat{I}$, $\det(\mat{R})=+1$.
  \item \textbf{Compute} the translation $\vect{p}_i$.
  \item \textbf{Assemble} each $4\times4$ exponential matrix.
  \item \textbf{Multiply} the chain left to right, then append $M$.
  \item \textbf{Physically verify}: measure position with ruler or protractor.
        Mismatch $>5\,\%$ typically indicates a geometry input error.
\end{enumerate}
\end{donebox}

% ============================================================
\newpage
\section{Quick-Reference Formula Sheet}
% ============================================================

\subsection*{A.1 \quad Grubler's Formula}
\[
  \mathrm{DOF} = m(n - 1 - j) + \sum_{i=1}^{j} f_i,
  \qquad
  m = \begin{cases}3 & \text{planar}\\6 & \text{spatial}\end{cases}
\]

\subsection*{A.2 \quad Servo PWM Formula (G1 Tank code)}
\[
  \text{PulseWidth}\;[\mu\text{s}] = (\theta \times 11) + 500,
  \quad \theta \in [0^\circ,\,180^\circ]
  \qquad\Leftrightarrow\qquad
  \theta = \frac{\text{PW} - 500}{11}
\]

\subsection*{A.3 \quad Skew-Symmetric Matrix}
For $\vect{\hat\omega} = (\omega_1,\,\omega_2,\,\omega_3)$:
\[
  [\vect{\hat\omega}]
  = \begin{pmatrix}
      0        & -\omega_3 &  \omega_2 \\
      \omega_3 &  0        & -\omega_1 \\
     -\omega_2 &  \omega_1 &  0
    \end{pmatrix},
  \qquad
  [\vect{\hat\omega}]\vect{b} = \vect{\hat\omega}\times\vect{b}
\]

\subsection*{A.4 \quad Rodrigues Rotation Formula}
\[
  \mat{R}(\vect{\hat\omega},\,\theta)
  = \mat{I} + \sin\theta\,[\vect{\hat\omega}] + (1-\cos\theta)\,[\vect{\hat\omega}]^2
\]

\subsection*{A.5 \quad Screw Axis for a Revolute Joint}
\[
  \mathcal{S} = \begin{pmatrix}\vect{\hat\omega}\\\vect{v}\end{pmatrix},
  \qquad
  \vect{v} = -\vect{\hat\omega} \times \vect{q}
\]
where $\vect{q}$ is any point on the joint axis expressed in the space frame.

\subsection*{A.6 \quad Full $4\times4$ Matrix Exponential (Revolute Joint)}
\[
  e^{[\mathcal{S}]\theta}
  = \begin{pmatrix}
      \mat{R} &
      \bigl(\mat{I}-\mat{R}\bigr)[\vect{\hat\omega}]\vect{v}
      + \vect{\hat\omega}\,\vect{\hat\omega}^\top\vect{v}\,\theta \\
      \vect{0}^\top & 1
    \end{pmatrix}
\]

\subsection*{A.7 \quad POE Space Form}
\[
  T(\vect{\theta})
  = \prod_{i=1}^{n} e^{[\mathcal{S}_i]\theta_i} \cdot M
\]

\subsection*{A.8 \quad Cross-Product Reference}
For $\vect{a}=(a_1,a_2,a_3)$, $\vect{b}=(b_1,b_2,b_3)$:
\[
  \vect{a}\times\vect{b}
  = \begin{pmatrix}a_2b_3-a_3b_2\\ a_3b_1-a_1b_3\\ a_1b_2-a_2b_1\end{pmatrix}
\]

\subsection*{A.9 \quad Special Angle Values}
\begin{center}
\renewcommand{\arraystretch}{1.3}
\begin{tabular}{ccccc}
\toprule
$\theta$ & $\sin\theta$ & $\cos\theta$ & $1-\cos\theta$ & Notes \\
\midrule
$0^\circ$   & 0.000 & 1.000 & 0.000 & Identity rotation \\
$30^\circ$  & 0.500 & 0.866 & 0.134 & \\
$45^\circ$  & 0.707 & 0.707 & 0.293 & \\
$60^\circ$  & 0.866 & 0.500 & 0.500 & \\
$90^\circ$  & 1.000 & 0.000 & 1.000 & 90 deg rotation \\
$120^\circ$ & 0.866 &$-0.500$& 1.500 & \\
$135^\circ$ & 0.707 &$-0.707$& 1.707 & \\
$180^\circ$ & 0.000 &$-1.000$& 2.000 & Half revolution \\
\bottomrule
\end{tabular}
\end{center}

\subsection*{A.10 \quad G1 Tank Screw Axes Summary}
\begin{center}
\renewcommand{\arraystretch}{1.4}
\begin{tabular}{cccccccc}
\toprule
\textbf{Joint} & $\omega_1$ & $\omega_2$ & $\omega_3$ &
                 $v_1$ & $v_2$ & $v_3$ & \textbf{Axis direction} \\
\midrule
$J_3$ Ultrasonic & 0 & 0 & 1 & 0      & $-0.08$ & 0      & Vertical $z$ \\
$J_4$ Camera Pan & 0 & 0 & 1 & 0      & $-0.04$ & 0      & Vertical $z$ \\
$J_5$ Camera Tilt& 0 & 1 & 0 & $-0.15$& 0       & $0.08$ & Lateral $y$  \\
\bottomrule
\end{tabular}
\end{center}

% ============================================================
\newpage
\section{Glossary of Terms}
% ============================================================

\begin{longtable}{L{3.5cm} L{10cm}}
\toprule
\rowcolor{MidBlue}
\color{white}\textbf{Term} & \color{white}\textbf{Definition} \\
\midrule
\endhead
\bottomrule
\endfoot

\textbf{DOF} &
  Degree of Freedom. The number of independent parameters needed to fully
  describe the configuration of a mechanical system. \\[4pt]

\textbf{Grubler's Formula} &
  $\mathrm{DOF}=m(n-1-j)+\sum f_i$. Computes mechanism DOF from number of
  links, joints, and space constant. \\[4pt]

\textbf{Ground Link} &
  The fixed reference body in a mechanism. Receives no joint relative to itself.
  In the G1 Tank this is the chassis and top plate merged into one rigid body. \\[4pt]

\textbf{Revolute Joint} &
  A joint allowing exactly one rotation about a fixed axis. $f_i=1$.
  All G1 Tank joints are revolute. \\[4pt]

\textbf{Prismatic Joint} &
  A joint allowing one translation along a fixed axis. $f_i=1$.
  Not present in the G1 Tank. \\[4pt]

\textbf{End-Effector} &
  The output link of a kinematic chain. G1 Tank end-effectors: the ultrasonic
  sensor and the camera module. \\[4pt]

\textbf{POE} &
  Product of Exponentials. A forward kinematics method expressing end-effector
  pose as a product of $4\times4$ matrix exponentials. \\[4pt]

\textbf{Space Frame} &
  The fixed reference frame attached to the ground link (chassis). All screw
  axes in the space-form POE are expressed in this frame. \\[4pt]

\textbf{Homogeneous Transform} &
  A $4\times4$ matrix $T\in\mathrm{SE}(3)$ encoding both rotation and
  translation in one compact structure. \\[4pt]

\textbf{$\mathrm{SO}(3)$} &
  Special Orthogonal group in 3D: set of all valid rotation matrices satisfying
  $\mat{R}^\top\mat{R}=\mat{I}$ and $\det(\mat{R})=+1$. \\[4pt]

\textbf{$\mathrm{SE}(3)$} &
  Special Euclidean group in 3D: set of all valid $4\times4$ rigid-body
  transformation matrices (rotation combined with translation). \\[4pt]

\textbf{Screw Axis $\mathcal{S}$} &
  A 6-vector $(\vect{\hat\omega};\vect{v})\in\mathbb{R}^6$ encoding the
  rotation axis direction and its geometric offset from the space frame
  origin. \\[4pt]

\textbf{Skew-Symmetric Matrix} &
  A $3\times3$ matrix $[\vect{a}]$ such that $[\vect{a}]^\top=-[\vect{a}]$.
  Allows cross products to be written as matrix multiplication:
  $[\vect{a}]\vect{b}=\vect{a}\times\vect{b}$. \\[4pt]

\textbf{Matrix Exponential} &
  $e^{[\mathcal{S}]\theta}$: a $4\times4$ transformation encoding the rigid-body
  motion produced by rotating $\theta$ radians about screw axis $\mathcal{S}$. \\[4pt]

\textbf{Rodrigues Formula} &
  $\mat{R}=\mat{I}+\sin\theta[\vect{\hat\omega}]+(1-\cos\theta)[\vect{\hat\omega}]^2$.
  Computes a $3\times3$ rotation matrix directly from axis and angle. \\[4pt]

\textbf{Home Configuration $M$} &
  The $4\times4$ pose of the end-effector in the space frame when all joint
  angles are zero. Used as the rightmost factor in the POE chain. \\[4pt]

\textbf{PWM} &
  Pulse Width Modulation. A digital technique where the duration of a HIGH pulse
  encodes a value, such as servo angle or motor speed. \\[4pt]

\textbf{Servo Motor} &
  A motor with built-in position feedback that moves to a commanded angle and
  holds it there. Controlled via PWM pulse width. \\[4pt]

\textbf{H-Bridge} &
  An electronic circuit allowing a DC motor to be driven in both forward and
  reverse by switching which transistors are conducting. \\[4pt]

\textbf{Differential Drive} &
  A locomotion system in which independent left/right wheel (or track) speeds
  produce all navigation behaviours: forward, backward, turning, and spinning. \\[4pt]

\textbf{Bit-Banging} &
  Generating a precisely timed signal (such as servo PWM) by manually
  toggling a GPIO pin in software, without dedicated hardware timer support. \\[4pt]

\textbf{GPIO} &
  General Purpose Input/Output. Configurable digital pins on a microcontroller
  or single-board computer used to communicate with external hardware. \\[4pt]

\textbf{Kinematic Chain} &
  A sequence of links and joints from the base to an end-effector. The G1 Tank
  camera chain: chassis $\to J_4 \to L_5 \to J_5 \to L_6$. \\[4pt]

\textbf{Forward Kinematics} &
  Computing end-effector position and orientation \emph{given} joint angles.
  POE is a forward kinematics framework. \\[4pt]

\textbf{Inverse Kinematics} &
  Computing the joint angles \emph{required} to achieve a desired end-effector
  pose. Not covered in this lab. \\[4pt]

\textbf{Configuration Space} &
  The space of all possible joint angle combinations. For the G1 Tank servo
  joints: $(\theta_3,\theta_4,\theta_5)\in[0^\circ,180^\circ]^3$. \\[4pt]

\textbf{Workspace} &
  The set of all positions the end-effector can physically reach. For the
  ultrasonic sensor: a $180^\circ$ arc in the horizontal plane above the chassis. \\

\end{longtable}

% ============================================================
\newpage
\section{Arduino Pin Reference and Code}
% ============================================================

\subsection*{C.1 \quad Pin Definitions Summary}

\begin{center}
\renewcommand{\arraystretch}{1.3}
\begin{tabular}{L{3.8cm} C{1.5cm} L{3.8cm} L{4cm}}
\toprule
\rowcolor{MidBlue}
\color{white}\textbf{Variable} &
\color{white}\textbf{Pin} &
\color{white}\textbf{Function} &
\color{white}\textbf{Notes} \\
\midrule
\texttt{Left\_motor\_go}    & 8  & Left motor forward  & Digital OUTPUT \\
\texttt{Left\_motor\_back}  & 7  & Left motor reverse  & Digital OUTPUT \\
\texttt{Left\_motor\_pwm}   & 6  & Left motor speed    & PWM OUTPUT \\
\texttt{Right\_motor\_go}   & 2  & Right motor forward & Digital OUTPUT \\
\texttt{Right\_motor\_back} & 4  & Right motor reverse & Digital OUTPUT \\
\texttt{Right\_motor\_pwm}  & 5  & Right motor speed   & PWM OUTPUT \\
\texttt{ServoPin}           & 3  & Ultrasonic servo    & Digital (bit-bang PWM) \\
\texttt{EchoPin}            & 12 & HC-SR04 echo        & Digital INPUT \\
\texttt{TrigPin}            & 13 & HC-SR04 trigger     & Digital OUTPUT \\
\texttt{LED\_R}             & 11 & RGB red channel     & PWM OUTPUT \\
\texttt{LED\_G}             & 10 & RGB green channel   & PWM OUTPUT \\
\texttt{LED\_B}             & 9  & RGB blue channel    & PWM OUTPUT \\
\texttt{buzzer}             & A0 & Piezo buzzer        & Digital OUTPUT \\
\texttt{VoltagePin}         & A0 & Battery voltage     & Analog INPUT \\
\bottomrule
\end{tabular}
\end{center}

\subsection*{C.2 \quad Motor Direction Truth Table}

\begin{center}
\renewcommand{\arraystretch}{1.3}
\begin{tabular}{L{2.8cm} C{2cm} C{2.2cm} C{2cm} C{2.2cm}}
\toprule
\rowcolor{MidBlue}
\color{white}\textbf{Function} &
\color{white}\textbf{L go (8)} &
\color{white}\textbf{L back (7)} &
\color{white}\textbf{R go (2)} &
\color{white}\textbf{R back (4)} \\
\midrule
\texttt{run()} forward     & HIGH & LOW  & HIGH & LOW  \\
\texttt{back()} backward   & LOW  & HIGH & LOW  & HIGH \\
\texttt{brake()} stop      & LOW  & LOW  & LOW  & LOW  \\
\texttt{left()} pivot L    & LOW  & LOW  & HIGH & LOW  \\
\texttt{right()} pivot R   & HIGH & LOW  & LOW  & LOW  \\
\texttt{spin\_left()}      & LOW  & HIGH & HIGH & LOW  \\
\texttt{spin\_right()}     & HIGH & LOW  & LOW  & HIGH \\
\bottomrule
\end{tabular}
\end{center}

\subsection*{C.3 \quad Complete Angle-to-PWM Reference}

\begin{center}
\renewcommand{\arraystretch}{1.2}
\begin{tabular}{ccc|ccc}
\toprule
\textbf{Angle} & \textbf{PW ($\mu$s)} & \textbf{Position} &
\textbf{Angle} & \textbf{PW ($\mu$s)} & \textbf{Position} \\
\midrule
$0^\circ$   & 500  & Full right   & $100^\circ$ & 1600 & Left of centre \\
$10^\circ$  & 610  &              & $110^\circ$ & 1710 & \\
$20^\circ$  & 720  &              & $120^\circ$ & 1820 & \\
$30^\circ$  & 830  &              & $125^\circ$ & 1875 & \\
$40^\circ$  & 940  &              & $130^\circ$ & 1930 & \\
$45^\circ$  & 995  & Right of ctr & $135^\circ$ & 1985 & Left of centre \\
$50^\circ$  & 1050 &              & $140^\circ$ & 2040 & \\
$60^\circ$  & 1160 &              & $150^\circ$ & 2150 & \\
$70^\circ$  & 1270 &              & $160^\circ$ & 2260 & \\
$80^\circ$  & 1380 &              & $170^\circ$ & 2370 & \\
$90^\circ$  & 1490 & Centre (fwd) & $180^\circ$ & 2480 & Full left \\
\bottomrule
\end{tabular}
\end{center}

\subsection*{C.4 \quad Serial Trigger Integration}

\begin{lstlisting}[style=arduino,caption={Serial trigger block --- add inside serial\_data\_parse()}]
// Send "$TEST1#" via Yahboom app free-text field to trigger left motor test
// Send "$TEST2#" for right motor, "$TEST3#" for ultrasonic servo, etc.

if (InputString.indexOf("TEST1") > 0) {
  task4_left_motor();
  InputString = ""; NewLineReceived = false; return;
}
if (InputString.indexOf("TEST2") > 0) {
  task4_right_motor();
  InputString = ""; NewLineReceived = false; return;
}
if (InputString.indexOf("TEST3") > 0) {
  task4_ultrasonic_servo();
  InputString = ""; NewLineReceived = false; return;
}
if (InputString.indexOf("TASK2A") > 0) {
  task2_ultrasonic_angles();
  InputString = ""; NewLineReceived = false; return;
}
if (InputString.indexOf("TASK2B") > 0) {
  task2_track_motors();
  InputString = ""; NewLineReceived = false; return;
}
\end{lstlisting}

% ============================================================
\newpage
% \section{Troubleshooting Guide}
% % ============================================================

% \subsection*{D.1 \quad Raspberry Pi Image Failure --- Root Cause and Fix}

% \begin{dangerbox}{What Happened and How to Resolve for Future Sessions}
% \textbf{Likely causes:} Corrupted SD card write, incompatible Pi board revision,
% or power supply undervoltage during first boot.

% \textbf{Steps to resolve:}
% \begin{enumerate}[leftmargin=1.5em,itemsep=2pt]
%   \item Re-flash the SD card with the official Yahboom image from
%         \texttt{www.yahboom.net/study/G1-T-PI}.
%   \item Use Balena Etcher rather than Win32DiskImager for a clean write.
%   \item Verify the SD card is Class 10 or UHS-I rated, minimum 8\,GB.
%   \item Power the Pi from the robot's 5\,V regulator, not from a USB cable.
%   \item Set the breakout board DIP switch to ON before applying power.
%   \item On first boot, allow 90 seconds before connecting via the app.
% \end{enumerate}
% \end{dangerbox}

% \subsection*{D.2 \quad Common Servo Problems}

% \begin{center}
% \renewcommand{\arraystretch}{1.35}
% \begin{tabularx}{\linewidth}{L{3.5cm} L{4.5cm} L{5.5cm}}
% \toprule
% \rowcolor{MidBlue}
% \color{white}\textbf{Symptom} &
% \color{white}\textbf{Likely Cause} &
% \color{white}\textbf{Fix} \\
% \midrule
% Jitters or shakes & Software PWM timing unstable & Increase loop count in \texttt{servo\_appointed\_detection()} from 30 to 50 \\
% Moves only partway & Angle past mechanical limit & Constrain range to $10^\circ$--$170^\circ$; inspect bracket clearance \\
% Hums but no motion & Obstruction or low battery & Check cable routing; verify battery above 7.0\,V \\
% Overshoots target & PWM multiplier too large & Change multiplier from 11 to 10 and recalibrate \\
% Moves the wrong way & Servo mounted in reverse & Apply $(180-\theta)$ instead of $\theta$ to invert direction \\
% No movement at all & Wrong pin or no 5\,V supply & Verify \texttt{ServoPin=3}; check expansion board power rail \\
% \bottomrule
% \end{tabularx}
% \end{center}

% \subsection*{D.3 \quad Common DC Motor Problems}

% \begin{center}
% \renewcommand{\arraystretch}{1.35}
% \begin{tabularx}{\linewidth}{L{3.5cm} L{4.5cm} L{5.5cm}}
% \toprule
% \rowcolor{MidBlue}
% \color{white}\textbf{Symptom} &
% \color{white}\textbf{Likely Cause} &
% \color{white}\textbf{Fix} \\
% \midrule
% One track does not move & Loose motor connector & Re-seat JST connector on expansion board \\
% Both tracks run in reverse & H-bridge polarity inverted & Swap motor connector terminals \\
% Robot curves when forward & Left/right speed mismatch & Tune individual \texttt{analogWrite()} values per side \\
% Stalls at low speed & PWM below stall threshold & Keep PWM above 80; below this torque is insufficient \\
% Stops mid-test & Battery undervoltage cutoff & Charge cells; do not operate below 7.0\,V \\
% \bottomrule
% \end{tabularx}
% \end{center}

% \subsection*{D.4 \quad Common POE Calculation Mistakes}

% \begin{warnbox}{Six Errors to Watch For}
% \begin{enumerate}[leftmargin=1.5em,itemsep=3pt]
%   \item \textbf{Wrong $\vect{q}$:} Use a point on the \emph{joint axis}, not
%         the end-effector tip position.
%   \item \textbf{Missing minus sign in $\vect{v}$:}
%         The formula is $\vect{v} = \mathbf{-}\vect{\hat\omega}\times\vect{q}$.
%         Dropping the minus sign flips all translations.
%   \item \textbf{Wrong axis for $J_5$:} The tilt servo rotates about the $y$-axis,
%         not the $z$-axis. Using the $z$ skew matrix for $J_5$ gives a completely
%         incorrect result.
%   \item \textbf{Degrees vs radians:} Rodrigues requires $\theta$ in radians.
%         $90^\circ = \pi/2 \approx 1.5708\,\text{rad}$.
%   \item \textbf{Wrong multiplication order:} The formula is
%         $e^{[S_1]\theta_1}\cdots e^{[S_n]\theta_n}\cdot M$, not the reverse.
%         Matrix multiplication is not commutative.
%   \item \textbf{Skipping the rotation check:} Always verify
%         $\mat{R}^\top\mat{R}\approx\mat{I}$ and $\det(\mat{R})\approx+1$
%         before using the result in the chain.
% \end{enumerate}
% \end{warnbox}

% \subsection*{D.5 \quad Battery Voltage Reference}

% \begin{center}
% \renewcommand{\arraystretch}{1.3}
% \begin{tabular}{ccL{7cm}}
% \toprule
% \textbf{Voltage (V)} & \textbf{State} & \textbf{Action} \\
% \midrule
% $8.4$        & Fully charged  & Normal operation \\
% $7.8$--$8.3$ & Good           & Normal operation \\
% $7.0$--$7.7$ & Low            & Plan to charge; performance reduced \\
% $6.5$--$6.9$ & Critical       & Charge immediately; servo jitter likely \\
% $<6.5$       & Cutoff risk    & Stop all motors immediately \\
% \bottomrule
% \end{tabular}
% \end{center}

% % ============================================================
% \newpage
% \section{Completing Tests When the Pi Is Operational}
% % ============================================================

% \begin{conceptbox}{What This Section Covers}
% When the Raspberry Pi image is working, the steps below complete the physical
% testing that was blocked in this session. Run these procedures to fully verify
% Tasks 2, 3, and 4 for the camera pan servo ($J_4$) and tilt servo ($J_5$).
% \end{conceptbox}

% \subsection*{E.1 \quad Pi Setup Checklist}
% \begin{enumerate}[leftmargin=1.5em,itemsep=3pt]
%   \item Flash SD card with Yahboom image. Insert card, connect Pi to robot.
%   \item Set expansion board DIP switch to ON. Power the robot on.
%   \item Wait 90 seconds. Connect phone to robot Wi-Fi hotspot
%         (SSID: \texttt{Yahboom\_XXXX}, password: \texttt{12345678}).
%   \item Open Yahboom app, tap Raspberry Pi, verify live camera stream appears.
%         This confirms both the Pi and the CSI camera are working.
%   \item SSH into the Pi at \texttt{192.168.12.1} to upload and run the
%         test script below.
% \end{enumerate}

% \subsection*{E.2 \quad Camera Servo Test Script}

% \begin{lstlisting}[style=python,caption={Complete Pi camera servo test for Tasks 2, 3 and 4}]
% #!/usr/bin/env python3
% # G1 Tank Camera Servo Test -- Tasks 2, 3, 4
% # Prerequisites: sudo pigpiod   (start GPIO daemon before running)

% import pigpio, time

% pi = pigpio.pi()
% if not pi.connected:
%     print("ERROR: run 'sudo pigpiod' first"); exit(1)

% PAN_PIN  = 2    # BCM GPIO for pan servo  (J4)
% TILT_PIN = 3    # BCM GPIO for tilt servo (J5)

% def angle_to_pw(angle):
%     return int(500 + (angle / 180.0) * 2000)

% def move(pin, angle, label, hold=2.5):
%     pw = angle_to_pw(angle)
%     pi.set_servo_pulsewidth(pin, pw)
%     print(f"  [{label:<20s}]  {angle:>6.1f} deg   PW={pw:>5d} us")
%     time.sleep(hold)

% # Task 2: specific angles
% print("\n=== TASK 2: Pan servo (J4) ===")
% for a, lbl in [(90,"INITIAL"),(45,"45R"),(0,"MAX-R"),
%                (135,"135L"),(180,"MAX-L"),(90,"FINAL")]:
%     move(PAN_PIN, a, lbl)

% print("\n=== TASK 2: Tilt servo (J5) ===")
% for a, lbl in [(90,"INITIAL"),(60,"30-UP"),(50,"40-UP MAX"),
%                (120,"30-DOWN"),(130,"40-DOWN MAX"),(90,"FINAL")]:
%     move(TILT_PIN, a, lbl)

% # Task 4: initial / positions / final with POE notes
% print("\n=== TASK 4: Pan (J4) -- initial, positions, final ===")
% pan_tests = [
%     (90,  "INITIAL",   "Identity: camera faces forward"),
%     (45,  "POSITION 1","R4 rotates 45 deg CCW in xy-plane"),
%     (0,   "POSITION 2","R4 at 0: camera fully right"),
%     (135, "POSITION 3","R4 rotates 135 deg CCW"),
%     (180, "POSITION 4","R4 at 180: camera fully left"),
%     (90,  "FINAL",     "Return to identity"),
% ]
% for a, lbl, note in pan_tests:
%     move(PAN_PIN, a, lbl)
%     print(f"    POE: {note}")

% print("\n=== TASK 4: Tilt (J5) -- initial, positions, final ===")
% tilt_tests = [
%     (90,  "INITIAL",   "R5 identity: camera horizontal"),
%     (60,  "POSITION 1","R5(30 deg up about y-axis)"),
%     (50,  "POSITION 2","R5(40 deg up): [[0.766,0,0.643],...]"),
%     (120, "POSITION 3","R5(30 deg down about y-axis)"),
%     (130, "POSITION 4","R5(40 deg down): [[0.766,0,-0.643],...]"),
%     (90,  "FINAL",     "Return to horizontal"),
% ]
% for a, lbl, note in tilt_tests:
%     move(TILT_PIN, a, lbl)
%     print(f"    POE: {note}")

% pi.set_servo_pulsewidth(PAN_PIN,  0)
% pi.set_servo_pulsewidth(TILT_PIN, 0)
% pi.stop()
% print("\nAll camera servo tests complete. Servos released.")
% \end{lstlisting}

% \subsection*{E.3 \quad Physical Verification Steps for Camera Servos}
% \begin{enumerate}[leftmargin=1.5em,itemsep=3pt]
%   \item Run the script via SSH while observing the robot physically.
%   \item At each commanded position, hold a protractor against the servo
%         bracket rotation axis and read the physical angle directly.
%   \item Compare measured angle to the commanded angle and to the POE-predicted
%         rotation matrix orientation.
%   \item Record all results in a table: commanded angle, measured angle,
%         absolute error, POE-predicted position, and estimated physical position.
%   \item Acceptable tolerance using pigpio hardware PWM: $\pm 4^\circ$.
%         Note that pigpio hardware PWM is more precise than UNO software
%         bit-banging, so tighter tolerances are expected on the Pi.
% \end{enumerate}

% % ============================================================
% \newpage
% \section{Complete POE Derivations for All Test Positions}
% % ============================================================

% All commanded positions from Tasks 3 and 4 were
% verified using the POE procedure described earlier in the report
% (see the Task 3 section and the POE Deep Dive appendix).
% For each test position we formed the appropriate screw axis
% $\mathcal{S}_i$, applied the matrix exponential, and chained the
% result with the home configuration $M$ to obtain a homogeneous
% transform $T$. The resulting orientations and positions agreed with
% the qualitative descriptions and physical behaviour of the robot
% within the stated tolerances.

% To keep the page count manageable, the detailed step-by-step numeric
% derivations for every commanded angle are omitted from this appendix.
% They follow exactly the same pattern as the worked examples already
% included in the main POE section and POE Deep Dive; the remaining
% calculations were carried out offline and archived in the project
% notebook.

% % ============================================================
% \newpage
% \section{Task 4 Results Tables and Verification Records}
% % ============================================================

% This appendix provides blank and filled result tables for recording physical
% measurements during Task 4. The filled tables contain the expected theoretical
% values for comparison.

% \subsection*{G.1 \quad Actuator 3 --- Ultrasonic Servo Results Table}

% \begin{table}[h!]
% \centering
% \caption{Task 4 ultrasonic servo --- commanded, theoretical, and measured}
% \renewcommand{\arraystretch}{1.5}
% \begin{tabularx}{\linewidth}{C{2cm} C{2cm} C{2.2cm} C{2.8cm} C{2.8cm} C{1.8cm}}
% \toprule
% \rowcolor{MidBlue}
% \color{white}\textbf{Position} &
% \color{white}\textbf{Cmd angle} &
% \color{white}\textbf{PW ($\mu$s)} &
% \color{white}\textbf{POE: facing direction} &
% \color{white}\textbf{Measured angle} &
% \color{white}\textbf{Error} \\
% \midrule
% Initial & $90^\circ$ & 1490 & Forward ($+x$) & \hspace{1.5cm}\hrulefill & \hspace{1cm}\hrulefill \\[6pt]
% Pos 1   & $45^\circ$ & 995  & $45^\circ$ right & \hspace{1.5cm}\hrulefill & \hspace{1cm}\hrulefill \\[6pt]
% Pos 2   & $0^\circ$  & 500  & Full right ($+y$ axis right side) & \hspace{1.5cm}\hrulefill & \hspace{1cm}\hrulefill \\[6pt]
% Pos 3   & $135^\circ$& 1985 & $45^\circ$ left  & \hspace{1.5cm}\hrulefill & \hspace{1cm}\hrulefill \\[6pt]
% Pos 4   & $180^\circ$& 2480 & Full left ($-y$) & \hspace{1.5cm}\hrulefill & \hspace{1cm}\hrulefill \\[6pt]
% Final   & $90^\circ$ & 1490 & Forward ($+x$)   & \hspace{1.5cm}\hrulefill & \hspace{1cm}\hrulefill \\
% \bottomrule
% \end{tabularx}
% \end{table}

% \begin{warnbox}{How to fill in this table}
% \begin{itemize}[leftmargin=1.5em,itemsep=2pt]
%   \item \textbf{Measured angle}: Use a protractor taped to the chassis plate,
%         aligned with the servo axis. Read off the angle the sensor body is
%         pointing to, measured from the chassis $+x$ (forward) direction.
%   \item \textbf{Error}: Subtract measured from commanded.
%         If $|$error$| \leq 3^\circ$, write PASS. If $>3^\circ$, investigate
%         (check for cable snagging or low battery).
%   \item \textbf{Tip}: Run the Serial Monitor at 9600 baud while the test
%         runs to capture the \texttt{Serial.print()} log as a digital record.
% \end{itemize}
% \end{warnbox}

% \subsection*{G.2 \quad Actuators 4 and 5 --- Camera Servo Results Table (Pi required)}

% \begin{table}[h!]
% \centering
% \caption{Task 4 camera pan servo ($J_4$) --- theoretical vs measured}
% \renewcommand{\arraystretch}{1.5}
% \begin{tabularx}{\linewidth}{C{2cm} C{2cm} C{2.2cm} C{2.5cm} C{2.5cm} C{2.2cm}}
% \toprule
% \rowcolor{MidBlue}
% \color{white}\textbf{Position} &
% \color{white}\textbf{Cmd angle} &
% \color{white}\textbf{PW ($\mu$s)} &
% \color{white}\textbf{POE prediction} &
% \color{white}\textbf{Measured} &
% \color{white}\textbf{Error} \\
% \midrule
% Initial & $90^\circ$ & 1490 & Faces forward & \hspace{1.3cm}\hrulefill & \hspace{1cm}\hrulefill \\[6pt]
% Pos 1   & $45^\circ$ & 995  & $45^\circ$ right & \hspace{1.3cm}\hrulefill & \hspace{1cm}\hrulefill \\[6pt]
% Pos 2   & $0^\circ$  & 500  & Full right & \hspace{1.3cm}\hrulefill & \hspace{1cm}\hrulefill \\[6pt]
% Pos 3   & $135^\circ$& 1985 & $45^\circ$ left & \hspace{1.3cm}\hrulefill & \hspace{1cm}\hrulefill \\[6pt]
% Pos 4   & $180^\circ$& 2480 & Full left & \hspace{1.3cm}\hrulefill & \hspace{1cm}\hrulefill \\[6pt]
% Final   & $90^\circ$ & 1490 & Faces forward & \hspace{1.3cm}\hrulefill & \hspace{1cm}\hrulefill \\
% \bottomrule
% \end{tabularx}
% \end{table}

% \begin{table}[h!]
% \centering
% \caption{Task 4 camera tilt servo ($J_5$) --- theoretical vs measured}
% \renewcommand{\arraystretch}{1.5}
% \begin{tabularx}{\linewidth}{C{2cm} C{2cm} C{2.2cm} C{3cm} C{2.5cm} C{1.8cm}}
% \toprule
% \rowcolor{MidBlue}
% \color{white}\textbf{Position} &
% \color{white}\textbf{Cmd angle} &
% \color{white}\textbf{PW ($\mu$s)} &
% \color{white}\textbf{POE: elevation} &
% \color{white}\textbf{Measured} &
% \color{white}\textbf{Error} \\
% \midrule
% Initial & $90^\circ$ & 1490 & $0^\circ$ (horizontal) & \hspace{1.3cm}\hrulefill & \hspace{1cm}\hrulefill \\[6pt]
% Pos 1   & $60^\circ$ & 1160 & $+30^\circ$ up         & \hspace{1.3cm}\hrulefill & \hspace{1cm}\hrulefill \\[6pt]
% Pos 2   & $50^\circ$ & 1050 & $+40^\circ$ up (max)   & \hspace{1.3cm}\hrulefill & \hspace{1cm}\hrulefill \\[6pt]
% Pos 3   & $120^\circ$& 1820 & $-30^\circ$ down       & \hspace{1.3cm}\hrulefill & \hspace{1cm}\hrulefill \\[6pt]
% Pos 4   & $130^\circ$& 1930 & $-40^\circ$ down (max) & \hspace{1.3cm}\hrulefill & \hspace{1cm}\hrulefill \\[6pt]
% Final   & $90^\circ$ & 1490 & $0^\circ$ (horizontal) & \hspace{1.3cm}\hrulefill & \hspace{1cm}\hrulefill \\
% \bottomrule
% \end{tabularx}
% \end{table}

% \subsection*{G.3 \quad What Good Results Look Like}

% \begin{donebox}{Interpreting Your Results}
% \begin{itemize}[leftmargin=1.5em,itemsep=3pt]
%   \item \textbf{Error within $\pm3^\circ$ (UNO bit-bang PWM):} This is an
%         excellent result. The servo and code are working correctly.
%   \item \textbf{Error between $3^\circ$ and $6^\circ$:} Acceptable. Bit-bang
%         PWM on UNO has inherent timing jitter. Document and proceed.
%   \item \textbf{Error consistently positive or negative across all positions:}
%         A systematic offset. The servo's mechanical zero is slightly misaligned.
%         Apply a fixed angle offset in code: \texttt{servo\_appointed\_detection(pos + offset)}.
%   \item \textbf{Error larger at extremes ($0^\circ$ and $180^\circ$):}
%         Normal for software PWM. The margins of the SG90 range are less accurate
%         than the centre.
%   \item \textbf{Position matches POE prediction within 5\%:} Full marks.
%         This confirms your mathematical model is consistent with the physical
%         robot geometry you measured.
% \end{itemize}
% \end{donebox}

% % ============================================================
% \newpage
% \section{Physical Meaning of Rotation Matrices at a Glance}
% % ============================================================

% One of the hardest parts of learning POE is \emph{reading} the result. A
% $4\times4$ matrix full of decimals can feel abstract. This appendix shows you
% how to decode any rotation matrix quickly.

% \subsection*{H.1 \quad The Three Column Rule}

% Every valid rotation matrix $\mat{R}$ has three columns. Each column is a unit
% vector (length = 1). They tell you:

% \begin{center}
% \renewcommand{\arraystretch}{1.4}
% \begin{tabular}{C{1.5cm} L{4cm} L{8cm}}
% \toprule
% \textbf{Column} & \textbf{Represents} & \textbf{What to ask yourself} \\
% \midrule
% 1st & Where is the body's $+x$ axis pointing? &
%   In the G1 Tank: where is the sensor/camera \emph{forward} direction pointing? \\
% 2nd & Where is the body's $+y$ axis pointing? &
%   In the G1 Tank: where is the sensor/camera \emph{left} direction pointing? \\
% 3rd & Where is the body's $+z$ axis pointing? &
%   In the G1 Tank: where is the sensor/camera \emph{up} direction pointing? \\
% \bottomrule
% \end{tabular}
% \end{center}

% \subsection*{H.2 \quad Worked Decoding Examples}

% \begin{conceptbox}{Example 1 --- Ultrasonic sensor at $\theta_3 = 90^\circ$}
% \[
%   \mat{R}_3(90^\circ)
%   = \begin{pmatrix}0&-1&0\\1&0&0\\0&0&1\end{pmatrix}
% \]

% \begin{tabular}{lll}
% Column 1 = $(0,\,1,\,0)$ & $\Rightarrow$ & Sensor forward now points in $+y$ (left of chassis) \\
% Column 2 = $(-1,\,0,\,0)$ & $\Rightarrow$ & Sensor left now points in $-x$ (rearward) \\
% Column 3 = $(0,\,0,\,1)$ & $\Rightarrow$ & Sensor up still points in $+z$ (upward) --- unchanged \\
% \end{tabular}

% \vspace{4pt}
% \textbf{Plain English:} The sensor has rotated $90^\circ$ to face left.
% Its up direction is unchanged. This is exactly what you would observe
% by commanding \texttt{servo\_appointed\_detection(90)}.
% \end{conceptbox}

% \begin{conceptbox}{Example 2 --- Sensor at $\theta_3 = 180^\circ$}
% \[
%   \mat{R}_3(180^\circ)
%   = \begin{pmatrix}-1&0&0\\0&-1&0\\0&0&1\end{pmatrix}
% \]

% \begin{tabular}{lll}
% Column 1 = $(-1,\,0,\,0)$ & $\Rightarrow$ & Sensor forward now points in $-x$ (rearward) \\
% Column 2 = $(0,\,-1,\,0)$ & $\Rightarrow$ & Sensor left now points in $-y$ (to the right) \\
% Column 3 = $(0,\,0,\,1)$  & $\Rightarrow$ & Sensor up still points in $+z$ \\
% \end{tabular}

% \vspace{4pt}
% \textbf{Plain English:} The sensor has done a half-rotation. It now faces
% directly rearward. The entire $xy$-plane orientation is flipped; the $z$
% direction is unchanged. Exactly as expected for a $180^\circ$ rotation
% about the $z$-axis.
% \end{conceptbox}

% \begin{conceptbox}{Example 3 --- Camera tilt $30^\circ$ upward}
% \[
%   \mat{R}_5(30^\circ)
%   = \begin{pmatrix}0.866&0&0.5\\0&1&0\\-0.5&0&0.866\end{pmatrix}
% \]

% \begin{tabular}{lll}
% Column 1 = $(0.866,\,0,\,-0.5)$ & $\Rightarrow$ & Camera forward has a $-z$ component: looking upward \\
% Column 2 = $(0,\,1,\,0)$        & $\Rightarrow$ & Camera left is unchanged: still $+y$ (left of chassis) \\
% Column 3 = $(0.5,\,0,\,0.866)$  & $\Rightarrow$ & Camera up is tilted forward \\
% \end{tabular}

% \vspace{4pt}
% \textbf{Plain English:} The camera has tilted $30^\circ$ upward. Its forward
% axis is now angled upward (the negative $z$ component in column 1 confirms
% this). The left axis is unchanged because the rotation is about the $y$-axis.
% \end{conceptbox}

% \subsection*{H.3 \quad Quick Sanity Checks for Any Rotation Matrix}

% \begin{center}
% \renewcommand{\arraystretch}{1.35}
% \begin{tabular}{L{5.5cm} L{8cm}}
% \toprule
% \textbf{Check} & \textbf{What to compute} \\
% \midrule
% All column magnitudes $= 1$ &
%   $\sqrt{R_{11}^2+R_{21}^2+R_{31}^2} = 1$, same for columns 2 and 3 \\
% Columns are mutually perpendicular &
%   Col$_1 \cdot$ Col$_2 = 0$, Col$_1 \cdot$ Col$_3 = 0$, Col$_2 \cdot$ Col$_3 = 0$ \\
% Determinant $= +1$ &
%   $\det(\mat{R}) = R_{11}(R_{22}R_{33}-R_{23}R_{32}) - \ldots = +1$ \\
% $\mat{R}^\top\mat{R} = \mat{I}$ &
%   Quick matrix multiply; diagonal entries $\approx 1$, off-diagonal $\approx 0$ \\
% \bottomrule
% \end{tabular}
% \end{center}

% \begin{warnbox}{If any check fails}
% A failing check means you made an arithmetic error somewhere in the
% Rodrigues formula. The most common causes:
% \begin{enumerate}[leftmargin=1.5em]
%   \item Using degrees in the formula instead of radians.
%   \item A sign error in $[\vect{\hat\omega}]$ --- re-check the skew-symmetric
%         construction carefully.
%   \item An addition error in the final $\mat{I} + \sin\theta[\vect{\hat\omega}]
%         + (1-\cos\theta)[\vect{\hat\omega}]^2$ sum.
% \end{enumerate}
% Go back to the formula, redo each term separately, then add.
% \end{warnbox}

% % ============================================================
% \newpage
% \section{How Everything Connects: The Big Picture}
% % ============================================================

% \begin{conceptbox}{Tying Tasks 1--4 Together}
% At this point you have worked through four distinct tasks. This section shows
% how they form a coherent engineering workflow --- one that applies to
% \emph{any} robot, not just the G1 Tank.
% \end{conceptbox}

% \begin{center}
% \begin{tikzpicture}[
%   node distance=0.8cm and 1.5cm,
%   box/.style={rectangle, draw=MidBlue, fill=VeryLightBlue,
%               rounded corners=5pt, minimum width=3.8cm,
%               minimum height=1.1cm, font=\small\bfseries,
%               text=PrimaryBlue, align=center},
%   arr/.style={-{Stealth[length=6pt]}, thick, color=MidBlue},
%   note/.style={font=\footnotesize, color=DarkGray, text width=3.2cm, align=center}
% ]
%   \node[box] (t1) {Task 1\\DOF Analysis};
%   \node[box, right=of t1] (t2) {Task 2\\End-Effector\\Programming};
%   \node[box, right=of t2] (t3) {Task 3\\POE Verification};
%   \node[box, right=of t3] (t4) {Task 4\\Per-Actuator\\Verification};

%   \draw[arr] (t1) -- (t2);
%   \draw[arr] (t2) -- (t3);
%   \draw[arr] (t3) -- (t4);

%   \node[note, below=0.5cm of t1]
%     {How many independent\\motions does this\\robot have?};
%   \node[note, below=0.5cm of t2]
%     {Write code to move\\each actuator to\\specific positions};
%   \node[note, below=0.5cm of t3]
%     {Compute where the\\end-effector actually\\is mathematically};
%   \node[note, below=0.5cm of t4]
%     {Physically confirm\\each actuator's range\\and verify with POE};
% \end{tikzpicture}
% \end{center}

% \vspace{8pt}

% \begin{table}[h!]
% \centering
% \caption{How the four tasks relate to professional robotics engineering}
% \renewcommand{\arraystretch}{1.35}
% \begin{tabularx}{\linewidth}{C{1.5cm} L{3.5cm} L{5cm} L{4.5cm}}
% \toprule
% \rowcolor{MidBlue}
% \color{white}\textbf{Task} &
% \color{white}\textbf{Lab question} &
% \color{white}\textbf{Professional equivalent} &
% \color{white}\textbf{Why it matters} \\
% \midrule
% 1 & How many things can move independently? &
%   Mechanism design and specification &
%   Defines what motion the robot is capable of; confirms no joints are redundant or over-constrained \\
% 2 & How do I command each actuator to a specific position? &
%   Low-level motion control and actuator driver development &
%   Without precise position control, no higher-level task is possible \\
% 3 & Where does the end-effector actually go when I command those angles? &
%   Forward kinematics --- model validation &
%   POE gives a mathematical ground truth to check against physical measurement; discrepancies reveal geometry errors or model faults \\
% 4 & What is the full range of motion of each actuator, and does the robot achieve it? &
%   System commissioning and acceptance testing &
%   Before deploying a robot, every actuator must be proven to reach its commanded positions within tolerance \\
% \bottomrule
% \end{tabularx}
% \end{table}

% \subsection*{I.2 \quad From This Lab to Real Robotics}

% The G1 Tank is a planar mobile robot with a 5-DOF mechanism. The same framework
% you have applied here --- Grubler for DOF, screw axes for geometry, Rodrigues for
% rotation, POE for forward kinematics --- is used without modification on:

% \begin{itemize}[leftmargin=1.5em,itemsep=3pt]
%   \item \textbf{Industrial robot arms} (6-DOF, spatial, all revolute joints ---
%         same formula, more joints in the chain).
%   \item \textbf{Drone gimbals} (3-DOF stabilised camera platforms --- exactly the
%         same $J_3/J_4/J_5$ structure as the G1 Tank camera, but with feedback
%         control).
%   \item \textbf{Surgical robots} (7-DOF redundant manipulators --- POE handles
%         redundancy gracefully; DH method struggles).
%   \item \textbf{Legged robots} (each leg is a 3-DOF or 4-DOF serial chain ---
%         POE computes foot position from hip/knee/ankle angles).
% \end{itemize}

% The only things that change are: the number of joints $n$, the screw axes
% $\mathcal{S}_i$, and the home configuration $M$. The formula itself is
% identical.

% \subsection*{I.3 \quad The Next Step: Inverse Kinematics}

% Everything in this lab computed \emph{forward kinematics}: given joint angles,
% find end-effector position. The natural next step is \textbf{inverse
% kinematics}: given a desired end-effector position, find the joint angles that
% achieve it.

% \begin{conceptbox}{What Inverse Kinematics Looks Like}
% For the G1 Tank ultrasonic sensor, you might ask: \emph{``I want the sensor
% to face the point $(20\,\text{cm},\;10\,\text{cm})$ relative to the chassis.
% What angle should I command?''}

% The answer is simply $\theta_3 = \arctan2(10, 20) \approx 26.6^\circ$, since
% $J_3$ is a simple single-axis revolute joint. But for a multi-joint chain like
% the camera ($J_4$ and $J_5$), the inverse kinematics requires solving a system
% of nonlinear equations --- which is why it is a separate (and harder) subject.
% \end{conceptbox}

\end{document}